\documentclass[8pt]{beamer}
\usetheme{Madrid}
\usecolortheme{default}
\usepackage{amsmath, amssymb, graphicx, array, multirow, booktabs, xcolor, tikz}
\usetikzlibrary{shapes,arrows,positioning,fit,backgrounds}

% Compact font settings
\setbeamerfont{title}{size=\Large}
\setbeamerfont{subtitle}{size=\small}
\setbeamerfont{author}{size=\scriptsize}
\setbeamerfont{date}{size=\scriptsize}
\setbeamerfont{frametitle}{size=\large}
\setbeamerfont{framesubtitle}{size=\normalsize}
\setbeamerfont{enumerate item}{size=\footnotesize}
\setbeamerfont{itemize item}{size=\footnotesize}
\setbeamerfont{block title}{size=\footnotesize}

% Compact spacing
\setlength{\itemsep}{0.08ex}
\setlength{\parskip}{0.08ex}
\setlength{\leftmargini}{0.3cm}
\setbeamersize{text margin left=3mm, text margin right=3mm}
\addtobeamertemplate{frame begin}{\vspace{-0.4cm}}{}

% Colors
\definecolor{myblue}{RGB}{0,0,255}
\definecolor{myred}{RGB}{255,0,0}
\definecolor{mygreen}{RGB}{0,128,0}
\definecolor{myorange}{RGB}{255,165,0}
\definecolor{mypurple}{RGB}{128,0,128}
\definecolor{mygray}{RGB}{128,128,128}
\definecolor{mygold}{RGB}{255,215,0}

\title{Data and AI Model Governance}
\subtitle{Complete Tutorial: Data Governance, Model Risk Management, Validation, GenAI Governance}
\author{GARP RAI Exam Preparation}
\date{}

\begin{document}
	
	\begin{frame}
		\titlepage
	\end{frame}
	
	% ============================================================
	% SECTION 1: INTRODUCTION TO DATA AND MODEL GOVERNANCE - SLIDES 1-5
	% ============================================================
	
	\begin{frame}{Slide 1: Introduction to Model Governance}
		\fontsize{6.5}{7.5}\selectfont
		\textbf{What is Model Governance?}
		
		\begin{itemize}
			\item Sound governance of models is essential for effective and prudent risk management
			\item Without it, a range of outcomes and decisions may occur that can prove detrimental
			\item Model governance should exist across the entire model lifecycle
			\item From model development through performance monitoring and ultimately decommissioning
		\end{itemize}
		
		\vspace{0.1cm}
		\textbf{Why Governance Matters for AI/ML:}
		\begin{itemize}
			\item Ensures technological advances don't overshadow need for fundamental challenges
			\item Addresses opacity of AI/ML models
			\item Proper governance of data used to train these models is particularly relevant
			\item Generative AI poses additional governance challenges
		\end{itemize}
		
		\vspace{0.1cm}
		\textbf{Key Insight:} Governance must span \textcolor{myblue}{the entire model lifecycle} from development to decommissioning!
		
		\textcolor{myblue}{\textbf{Focus:}} Financial sector and quantitative risk models (QRMs) due to established regulatory framework.
	\end{frame}
	
	\begin{frame}{Slide 2: Data Governance Definition}
		\fontsize{6.5}{7.5}\selectfont
		\textbf{What is Data Governance?}
		
		\vspace{0.1cm}
		\textbf{Data Governance Institute Definition:}
		\begin{itemize}
			\item "A system of decision rights and accountabilities for information-related processes"
			\item "Executed according to agreed-upon models which describe who can take what actions"
			\item "With what information, and when, under what circumstances, using what methods"
		\end{itemize}
		
		\vspace{0.1cm}
		\textbf{Broad Scope:}
		\begin{itemize}
			\item Encompasses the people, processes, and technologies required to manage and protect data assets
			\item Covers entire data lifecycle from collection to deletion
			\item Balances data accessibility with data protection
		\end{itemize}
		
		\vspace{0.1cm}
		\textbf{Key Components:}
		\begin{itemize}
			\item Data strategy and vision
			\item Data quality management
			\item Data provenance and legal compliance
			\item Data classification and security
			\item Metadata management
		\end{itemize}
		
		\textcolor{myblue}{\textbf{Key Insight:}} Data governance is about \textcolor{myblue}{decision rights and accountabilities} for data!
	\end{frame}
	
	\begin{frame}{Slide 3: Data Strategy - Vision and Goals}
		\fontsize{6.5}{7.5}\selectfont
		\textbf{Developing a Data Strategy:}
		
		\vspace{0.1cm}
		\textbf{Key Requirements:}
		\begin{itemize}
			\item Clear and approved data-governance policy
			\item Strategic vision and operational practices
			\item Level of authorization for non-publicly available personal information
			\item Specification of data usage (single or recurring use)
			\item Documentation of data sources
		\end{itemize}
		
		\vspace{0.1cm}
		\textbf{Alternative Data Considerations:}
		\begin{itemize}
			\item Additional controls and testing required
			\item Privacy concerns
			\item Data quality and accuracy risks
			\item Unbounded data streams
		\end{itemize}
		
		\vspace{0.1cm}
		\textbf{Data Governance Committee Responsibilities:}
		\begin{itemize}
			\item Data collection, reconciliations, treatment, access, monitoring
			\item Management of data stewards
			\item Overall oversight on all aspects of data governance
		\end{itemize}
		
		\textcolor{myblue}{\textbf{Key Insight:}} Data strategy requires \textcolor{myblue}{clear policies and authorization} for data use!
	\end{frame}
	
	\begin{frame}{4: Data Quality - Accuracy, Consistency, Integrity}
		\fontsize{6.5}{7.5}\selectfont
		\textbf{Ensuring Data Quality:}
		
		\vspace{0.1cm}
		\textbf{Core Requirements:}
		\begin{itemize}
			\item Complete, relevant, transparent, and accurate data should be the primary objective
			\item Model risk management requirements around data quality do not change with AI/ML algorithms
		\end{itemize}
		
		\vspace{0.1cm}
		\textbf{Alternative Data Sources - Additional Scrutiny:}
		\begin{itemize}
			\item Verifying accuracy of alternative data sources may be difficult
			\item Data providers may change field definitions after deployment
			\item Even well-established data sources are subject to this risk
		\end{itemize}
		
		\vspace{0.1cm}
		\textbf{Example - COVID-19 Pandemic:}
		\begin{itemize}
			\item U.S. Bureau of Labor Statistics unemployment data was biased
			\item Inability to conduct proper surveys
			\item Anomalies may occur with any data stream
			\item New data sources introduce additional risks
		\end{itemize}
		
		\vspace{0.1cm}
		\textbf{Data Monitoring:}
		\begin{itemize}
			\item Input data monitoring to alert model owners of flawed inputs
			\item Outlier detection may be insufficient to spot anomalies
			\item Definitional changes may only be disclosed in footnotes
		\end{itemize}
		
		\textcolor{myblue}{\textbf{Key Insight:}} Data quality is \textcolor{myblue}{critical and requires continuous monitoring}!
	\end{frame}
	
	\begin{frame}{Slide 5: Data Provenance}
		\fontsize{6.5}{7.5}\selectfont
		\textbf{Data Provenance - Verifying Legal Right to Use Data:}
		
		\vspace{0.1cm}
		\textbf{Definition:}
		\begin{itemize}
			\item Provenance refers to the sequence of ownership and handling of items of value
			\item Now applies to data used to create models
			\item Proving legal right to use data is essential
		\end{itemize}
		
		\vspace{0.1cm}
		\textbf{Why It Matters:}
		\begin{itemize}
			\item Scrutiny of whether vendors have legal right to use data
			\item Copyright and intellectual property lawsuits against GenAI vendors
			\item Example: Data scraped from websites without permission
		\end{itemize}
		
		\vspace{0.1cm}
		\textbf{Regulatory Consequences:}
		\begin{itemize}
			\item FTC: Any model built on improperly obtained data must be destroyed
			\item First such legal action has already been settled with model destruction
			\item No model can be approved without proving ownership of input data
			\item Vendor models need to provide proof and indemnify licensees
		\end{itemize}
		
		\vspace{0.1cm}
		\textbf{Key Challenge:}
		\begin{itemize}
			\item Established data (e.g., FICO) is generally safe
			\item Alternative data requires scrutiny
			\item Risk must be considered for GenAI models where model risk management applies
		\end{itemize}
		
		\textcolor{myblue}{\textbf{Key Insight:}} Data provenance is \textcolor{myblue}{essential for legal compliance and model approval}!
	\end{frame}
	
	% ============================================================
	% SECTION 2: DATA GOVERNANCE DETAILS - SLIDES 6-15
	% ============================================================
	
	\begin{frame}{Slide 6: Data Classification}
		\fontsize{6.5}{7.5}\selectfont
		\textbf{Data Classification - Structure and Confidentiality:}
		
		\vspace{0.1cm}
		\textbf{Data Types:}
		\begin{enumerate}
			\item \textbf{Quantitative Data:}
			\begin{itemize}
				\item Measured or determined numerically
				\item Examples: income, age, temperature
			\end{itemize}
			\item \textbf{Qualitative Data:}
			\begin{itemize}
				\item Determined non-numerically
				\item Examples: gender, ethnicity, job title
			\end{itemize}
		\end{enumerate}
		
		\vspace{0.1cm}
		\textbf{Data Sources:}
		\begin{itemize}
			\item Observed data
			\item Projected data
			\item Judgment-based data
		\end{itemize}
		
		\vspace{0.1cm}
		\textbf{Critical Governance Requirement:}
		\begin{itemize}
			\item Distinguish the classification of data being used
			\item Ensure access and use is permissible
			\item When data is confidential or personally identifiable, controls must be:
			\begin{itemize}
				\item Present
				\item Operational
				\item Effective
			\end{itemize}
		\end{itemize}
		
		\textcolor{myblue}{\textbf{Key Insight:}} Data classification determines \textcolor{myblue}{access controls and usage permissions}!
	\end{frame}
	
	\begin{frame}{Slide 7: Metadata Management}
		\fontsize{6.5}{7.5}\selectfont
		\textbf{Metadata Management:}
		
		\vspace{0.1cm}
		\textbf{What is Metadata?}
		\begin{itemize}
			\item Basic information about data
			\item Includes: file type, time of creation, data author
			\item Data source, file size, and other relevant information
		\end{itemize}
		
		\vspace{0.1cm}
		\textbf{Types of Metadata:}
		\begin{enumerate}
			\item \textbf{Descriptive:} Identifies and describes data
			\item \textbf{Structural:} How data is organized
			\item \textbf{Administrative:} Technical and rights information
			\item \textbf{Reference:} Standards and definitions
			\item \textbf{Statistical:} Data quality and methodology
		\end{enumerate}
		
		\vspace{0.1cm}
		\textbf{Robust Metadata Management Strategy:}
		\begin{itemize}
			\item Data documentation
			\item Data mapping
			\item Data dictionary
			\item Data definitions
			\item Data process flows
			\item Data relationships to other data
			\item Data structures
		\end{itemize}
		
		\textcolor{myblue}{\textbf{Key Insight:}} Metadata management enables \textcolor{myblue}{better-informed business decisions}!
	\end{frame}
	
	\begin{frame}{Slide 8: Data Security and Protection}
		\fontsize{6.5}{7.5}\selectfont
		\textbf{Data Security and Protection:}
		
		\vspace{0.1cm}
		\textbf{Key Concepts:}
		\begin{itemize}
			\item Data security often used interchangeably with data protection and data integrity
			\item Strategic security process and data protection plan are essential
			\item Procedural steps identified to safeguard:
			\begin{itemize}
				\item Availability of data
				\item Access to data
				\item Data privacy
			\end{itemize}
		\end{itemize}
		
		\vspace{0.1cm}
		\textbf{Data Protection Strategy:}
		\begin{itemize}
			\item Safeguarding important information from:
			\begin{itemize}
				\item Corruption
				\item Malicious or accidental damage
				\item Compromise or loss
			\end{itemize}
			\item Importance increases with amount of data created and stored
			\item Data retention policy should be in place and adhered to
		\end{itemize}
		
		\vspace{0.1cm}
		\textbf{US Regulatory Examples:}
		\begin{itemize}
			\item Safeguards Rule (FTC): Financial institutions must keep customer information secure
			\item Gramm-Leach-Bliley Act (GLBA): Privacy practices and opt-out rights
			\item Privacy Act of 1974: Federal agency data collection and use
		\end{itemize}
		
		\textcolor{myblue}{\textbf{Key Insight:}} Data security requires \textcolor{myblue}{comprehensive protection strategies and regulatory compliance}!
	\end{frame}
	
	\begin{frame}{Slide 9: Data Access and Sharing}
		\fontsize{6.5}{7.5}\selectfont
		\textbf{Data Access: Permissions and Secure Sharing:}
		
		\vspace{0.1cm}
		\textbf{Most Important Aspect:}
		\begin{itemize}
			\item Ensuring data is readily accessible and shareable
			\item To those who have requisite permission and need to access it
		\end{itemize}
		
		\vspace{0.1cm}
		\textbf{Components of Data Access:}
		\begin{itemize}
			\item Data security
			\item Data protection
			\item Data permissions
			\item Data access controls
		\end{itemize}
		
		\vspace{0.1cm}
		\textbf{Regulatory Requirements:}
		\begin{itemize}
			\item Organizations must implement appropriate security measures
			\item Encryption to safeguard data
			\item Ensure confidentiality and integrity
			\item Maintain compliance
		\end{itemize}
		
		\textcolor{myblue}{\textbf{Key Insight:}} Data access balances \textcolor{myblue}{availability with security and permissions}!
	\end{frame}
	
	\begin{frame}{Slide 10: Data Compliance}
		\fontsize{6.5}{7.5}\selectfont
		\textbf{Compliance: Legal and Regulatory Requirements:}
		
		\vspace{0.1cm}
		\textbf{Three Basic Principles:}
		\begin{enumerate}
			\item \textbf{Obtaining Consent:}
			\begin{itemize}
				\item Data subjects must consent to data collection
				\item Informed consent is essential
				\item Consent can be withdrawn
			\end{itemize}
			\item \textbf{Data Minimization:}
			\begin{itemize}
				\item Collect only necessary data
				\item Most effective privacy protection
				\item Reduce risk and exposure
			\end{itemize}
			\item \textbf{Ensuring Rights of Data Subjects:}
			\begin{itemize}
				\item Right to access data
				\item Right to correct data
				\item Right to delete data
				\item Right to be forgotten
			\end{itemize}
		\end{enumerate}
		
		\vspace{0.1cm}
		\textbf{Key Regulatory Frameworks:}
		\begin{itemize}
			\item GDPR (Europe)
			\item CCPA (California)
			\item PIPL (China)
			\item GLBA (US Financial)
		\end{itemize}
		
		\textcolor{myblue}{\textbf{Key Insight:}} Data compliance centers on \textcolor{myblue}{consent, minimization, and data subject rights}!
	\end{frame}
	
	\begin{frame}{Slide 11: Board Responsibilities in Data Governance}
		\fontsize{6.5}{7.5}\selectfont
		\textbf{Board of Directors' Role in Data Governance:}
		
		\vspace{0.1cm}
		\textbf{Crucial Responsibilities:}
		\begin{enumerate}
			\item \textbf{Approval of Framework:}
			\begin{itemize}
				\item Overall data governance framework and policies
				\item Clear guidelines on data quality, privacy, security
				\item Compliance with relevant regulations
			\end{itemize}
			\item \textbf{Oversight of Compliance and Risk Management:}
			\begin{itemize}
				\item Compliance with legal, regulatory, and ethical standards
				\item Data protection laws (GDPR), industry standards, internal policies
			\end{itemize}
			\item \textbf{Risk Assessment and Management:}
			\begin{itemize}
				\item Data breaches
				\item Data quality issues
				\item Misuse of data
			\end{itemize}
			\item \textbf{Regular Review:}
			\begin{itemize}
				\item Adapt to changing business needs
				\item New technologies
				\item Regulatory requirements
			\end{itemize}
		\end{enumerate}
		
		\textcolor{myblue}{\textbf{Key Insight:}} Board sets the \textcolor{myblue}{"tone at the top"} for data governance culture!
	\end{frame}
	
	\begin{frame}{Slide 12: Synthetic Data and AI Readiness}
		\fontsize{6.5}{7.5}\selectfont
		\textbf{Synthetic Data and AI Readiness:}
		
		\vspace{0.1cm}
		\textbf{Synthetic Data Usage:}
		\begin{itemize}
			\item Increasing use to augment available data
			\item Mitigates data risks (client or personally identifiable data)
			\item Can be designed to avoid biases
		\end{itemize}
		
		\vspace{0.1cm}
		\textbf{Advantages of Synthetic Data:}
		\begin{itemize}
			\item No privacy risks (doesn't come from real data subjects)
			\item Can be more easily rid of problematic biases by design
			\item Scalable and customizable
		\end{itemize}
		
		\vspace{0.1cm}
		\textbf{Disadvantages:}
		\begin{itemize}
			\item May not be as precise as real data (especially visual data)
			\item May differ from real data in ways not obvious
			\item Requires careful validation
		\end{itemize}
		
		\vspace{0.1cm}
		\textbf{AI-Ready Data:}
		\begin{itemize}
			\item Accurate, unbiased, well-structured
			\item Cleaned and suitable for training
			\item Improves accuracy and efficiency of AI models
		\end{itemize}
		
		\textcolor{myblue}{\textbf{Key Insight:}} Synthetic data and AI readiness \textcolor{myblue}{require consideration in data strategy}!
	\end{frame}
	
	\begin{frame}{Slide 13: Data Governance Process}
		\fontsize{6.5}{7.5}\selectfont
		\textbf{The Data Governance Process:}
		
		\vspace{0.1cm}
		\textbf{Key Steps:}
		\begin{enumerate}
			\item \textbf{Data Strategy Development:}
			\begin{itemize}
				\item Vision and goals
				\item Policies and procedures
				\item Authorization levels
			\end{itemize}
			\item \textbf{Data Collection:}
			\begin{itemize}
				\item Identify sources
				\item Verify provenance
				\item Ensure legal compliance
			\end{itemize}
			\item \textbf{Data Quality Management:}
			\begin{itemize}
				\item Accuracy checks
				\item Consistency validation
				\item Integrity verification
			\end{itemize}
			\item \textbf{Data Classification:}
			\begin{itemize}
				\item Structure and confidentiality
				\item Access controls
			\end{itemize}
			\item \textbf{Metadata Management:}
			\begin{itemize}
				\item Documentation
				\item Data dictionary
				\item Data mapping
			\end{itemize}
			\item \textbf{Data Security and Compliance:}
			\begin{itemize}
				\item Protection measures
				\item Regulatory compliance
				\item Monitoring and auditing
			\end{itemize}
		\end{enumerate}
		
		\textcolor{myblue}{\textbf{Key Insight:}} Data governance is a \textcolor{myblue}{continuous process across the data lifecycle}!
	\end{frame}
	
	\begin{frame}{Slide 14: Roles and Responsibilities in Data Governance}
		\fontsize{6.5}{7.5}\selectfont
		\textbf{Key Roles in Data Governance:}
		
		\vspace{0.1cm}
		\textbf{Board of Directors:}
		\begin{itemize}
			\item Approve framework and policies
			\item Oversee compliance and risk management
			\item Ensure regular review and updates
		\end{itemize}
		
		\vspace{0.1cm}
		\textbf{Data Governance Committee:}
		\begin{itemize}
			\item Overall oversight on all aspects
			\item Approve data practices and policies
			\item Not responsible for operational aspects
		\end{itemize}
		
		\vspace{0.1cm}
		\textbf{Data Stewards:}
		\begin{itemize}
			\item Day-to-day data management
			\item Ensure data quality and consistency
			\item Implement policies
		\end{itemize}
		
		\vspace{0.1cm}
		\textbf{Data Owners:}
		\begin{itemize}
			\item Accountability for specific data assets
			\item Define access permissions
			\item Ensure proper use
		\end{itemize}
		
		\textcolor{myblue}{\textbf{Key Insight:}} Clear roles and responsibilities are \textcolor{myblue}{essential for effective data governance}!
	\end{frame}
	
	\begin{frame}{Slide 15: Data Governance Summary}
		\fontsize{6.5}{7.5}\selectfont
		\textbf{Data Governance: Key Takeaways}
		
		\vspace{0.1cm}
		\textbf{Core Components:}
		\begin{enumerate}
			\item \textbf{Strategy:}
			\begin{itemize}
				\item Vision and goals
				\item Policies and authorization
				\item Alternative data controls
			\end{itemize}
			\item \textbf{Quality:}
			\begin{itemize}
				\item Accuracy, consistency, integrity
				\item Continuous monitoring
				\item Alternative data scrutiny
			\end{itemize}
			\item \textbf{Provenance:}
			\begin{itemize}
				\item Legal right to use data
				\item Copyright compliance
				\item Vendor verification
			\end{itemize}
			\item \textbf{Classification:}
			\begin{itemize}
				\item Structure and confidentiality
				\item Access controls
			\end{itemize}
			\item \textbf{Security:}
			\begin{itemize}
				\item Protection measures
				\item Regulatory compliance
				\item Data retention
			\end{itemize}
		\end{enumerate}
		
		\textcolor{myblue}{\textbf{Key Insight:}} Data governance is the \textcolor{myblue}{foundation for responsible AI}!
	\end{frame}
	
	% ============================================================
	% SECTION 3: MODEL GOVERNANCE - SLIDES 16-30
	% ============================================================
	
	\begin{frame}{Slide 16: What is a Model?}
		\fontsize{6.5}{7.5}\selectfont
		\textbf{Defining a Model (SR 11-7):}
		
		\vspace{0.1cm}
		\textbf{SR 11-7 Definition:}
		\begin{itemize}
			\item "The term 'model' refers to a quantitative method, system, or approach"
			\item "That applies statistical, economic, financial, or mathematical theories"
			\item "Techniques, and assumptions to process input data into quantitative estimates"
		\end{itemize}
		
		\vspace{0.1cm}
		\textbf{Key Elements:}
		\begin{enumerate}
			\item \textbf{Quantity of Interest:}
			\begin{itemize}
				\item Numerical object whose future value is uncertain
				\item Examples: portfolio value, revenue projection, client count
			\end{itemize}
			\item \textbf{Potential Future Scenarios:}
			\begin{itemize}
				\item Possible values of the quantity of interest
				\item Each scenario assigned a weight (probability)
			\end{itemize}
			\item \textbf{Risk Measure:}
			\begin{itemize}
				\item Summarizes essential information
				\item Examples: VaR, expected shortfall, average outcome
			\end{itemize}
		\end{enumerate}
		
		\textcolor{myblue}{\textbf{Key Insight:}} A model is a \textcolor{myblue}{quantitative method for processing data into estimates}!
	\end{frame}
	
	\begin{frame}{Slide 17: Risk Modeling Elements}
		\fontsize{6.5}{7.5}\selectfont
		\textbf{Three Main Elements of Quantitative Risk Modeling:}
		
		\vspace{0.1cm}
		\textbf{1. Quantity of Interest:}
		\begin{itemize}
			\item Numerical object whose future value is uncertain
			\item Examples:
			\begin{itemize}
				\item Portfolio value in ten business days
				\item Revenue projection over next five years
				\item Expected number of new clients by year-end
			\end{itemize}
		\end{itemize}
		
		\vspace{0.1cm}
		\textbf{2. Potential Future Scenarios:}
		\begin{itemize}
			\item Possible values of the quantity of interest
			\item Each scenario assigned a weight (probability)
			\item Examples:
			\begin{itemize}
				\item Portfolio value conditioned on investment decisions
				\item Revenue after business model changes
				\item New clients after advertising campaign
			\end{itemize}
		\end{itemize}
		
		\vspace{0.1cm}
		\textbf{3. Risk Measure:}
		\begin{itemize}
			\item Summarizes essential information from scenarios
			\item Examples:
			\begin{itemize}
				\item Value at Risk (VaR)
				\item Expected Shortfall (ES)
				\item Average outcome across all scenarios
			\end{itemize}
		\end{itemize}
		
		\textcolor{myblue}{\textbf{Key Insight:}} Risk models provide \textcolor{myblue}{structured approach to envision the future}!
	\end{frame}
	
	\begin{frame}{Slide 18: Important Model Considerations}
		\fontsize{6.5}{7.5}\selectfont
		\textbf{Key Considerations in Risk Modeling:}
		
		\vspace{0.1cm}
		\textbf{1. Completeness of Scenario Sets:}
		\begin{itemize}
			\item Challenging to anticipate every potential future scenario
			\item Historical data may not reflect rare events
			\item Expert judgment can be difficult
		\end{itemize}
		
		\vspace{0.1cm}
		\textbf{2. Feedback Effects:}
		\begin{itemize}
			\item Scenarios and decisions may influence other market participants
			\item Scenario sets and weights may need ongoing updates
		\end{itemize}
		
		\vspace{0.1cm}
		\textbf{3. Communication of Results:}
		\begin{itemize}
			\item Reports should provide summary of main assumptions
			\item Ensure transparency and avoid complacency
			\item Reflect perceived risk based on uncertainty and exposure
		\end{itemize}
		
		\vspace{0.1cm}
		\textbf{4. Regulatory Context:}
		\begin{itemize}
			\item SR 11-7 (2011) established regulatory standards
			\item OCC Handbook (2021) includes AI/ML perspectives
		\end{itemize}
		
		\textcolor{myblue}{\textbf{Key Insight:}} Effective risk modeling requires \textcolor{myblue}{completeness, awareness of feedback, and clear communication}!
	\end{frame}
	
	\begin{frame}{Slide 19: Model Governance Policies}
		\fontsize{6.5}{7.5}\selectfont
		\textbf{Model Governance Policies - Key Elements:}
		
		\vspace{0.1cm}
		\textbf{Core Policy Areas:}
		\begin{enumerate}
			\item \textbf{Defining What Constitutes a Model:}
			\begin{itemize}
				\item Distinction between model and end-user tool
				\item Gray areas need consistent adherence
				\item Presence of uncertainty may indicate a model
			\end{itemize}
			\item \textbf{Model Risk Definition:}
			\begin{itemize}
				\item Agree on internal definition
				\item Align with regulatory requirements
			\end{itemize}
			\item \textbf{Risk Tier Identification:}
			\begin{itemize}
				\item Drives depth and frequency of validation
				\item Materiality, complexity, exposure
			\end{itemize}
			\item \textbf{Model Validation Definition:}
			\begin{itemize}
				\item Establish expectations for validation process
				\item Documentation standards
			\end{itemize}
			\item \textbf{Model Risk Appetite:}
			\begin{itemize}
				\item Document institution's risk appetite
				\item Outline procedures for unacceptable risks
			\end{itemize}
		\end{enumerate}
		
		\textcolor{myblue}{\textbf{Key Insight:}} Written policies are the \textcolor{myblue}{foundation of model governance frameworks}!
	\end{frame}
	
	\begin{frame}{Slide 20: Model Documentation Standards}
		\fontsize{6.5}{7.5}\selectfont
		\textbf{Model Documentation - SR 11-7 Guidance:}
		
		\vspace{0.1cm}
		\textbf{Key Requirements:}
		\begin{itemize}
			\item Strong governance includes documentation of model development and validation
			\item Sufficiently detailed to allow parties unfamiliar with a model to understand:
			\begin{itemize}
				\item How the model operates
				\item Limitations of the model
				\item Key assumptions of the model
			\end{itemize}
		\end{itemize}
		
		\vspace{0.1cm}
		\textbf{Documentation Best Practices:}
		\begin{enumerate}
			\item \textbf{Comprehensive:}
			\begin{itemize}
				\item Cover all aspects of model lifecycle
				\item Include development and validation
				\item Document all key decisions
			\end{itemize}
			\item \textbf{Accessible:}
			\begin{itemize}
				\item Clear language
				\item Well-organized
				\item Executive summaries
			\end{itemize}
			\item \textbf{Maintained:}
			\begin{itemize}
				\item Updated with changes
				\item Version control
				\item Historical records
			\end{itemize}
		\end{enumerate}
		
		\textcolor{myblue}{\textbf{Key Insight:}} Documentation must be \textcolor{myblue}{sufficiently detailed for others to understand the model}!
	\end{frame}
	
	\begin{frame}{Slide 21: Model Inventory}
		\fontsize{6.5}{7.5}\selectfont
		\textbf{Model Inventory - Comprehensive Tracking:}
		
		\vspace{0.1cm}
		\textbf{What to Include:}
		\begin{enumerate}
			\item \textbf{Model Identification:}
			\begin{itemize}
				\item Model owner, user, and developer information
				\item Purpose and brief description
				\item Applicable products
			\end{itemize}
			\item \textbf{Methodology:}
			\begin{itemize}
				\item Description of model methodology
				\item Model classification (e.g., Monte Carlo, logistic regression)
				\item Whether it's an AI/ML model
			\end{itemize}
			\item \textbf{Development:}
			\begin{itemize}
				\item In-house or vendor developed
				\item Use and frequency
				\item Restrictions and limitations
			\end{itemize}
			\item \textbf{Assumptions and Limitations:}
			\begin{itemize}
				\item Key assumptions
				\item Known weaknesses
				\item Management overlays
			\end{itemize}
			\item \textbf{Documentation and History:}
			\begin{itemize}
				\item Locations of documentation and validation reports
				\item Access to model code and databases
				\item History of updates, approvals, validations
			\end{itemize}
		\end{enumerate}
		
		\textcolor{myblue}{\textbf{Key Insight:}} Model inventory enhances \textcolor{myblue}{transparency and tracking} of all models!
	\end{frame}
	
	\begin{frame}{Slide 22: Model Landscape and Interdependencies}
		\fontsize{6.5}{7.5}\selectfont
		\textbf{Model Landscape - Understanding Interdependencies:}
		
		\vspace{0.1cm}
		\textbf{Why Model Landscape Matters:}
		\begin{itemize}
			\item Models should not be viewed in isolation
			\item Outputs from one model often serve as inputs for others
			\item Transparency regarding interdependencies is vital
		\end{itemize}
		
		\vspace{0.1cm}
		\textbf{Key Considerations:}
		\begin{enumerate}
			\item \textbf{Interdependencies:}
			\begin{itemize}
				\item Upstream and downstream dependencies
				\item Impact on key business decisions
				\item Critical models identification
			\end{itemize}
			\item \textbf{Aggregate Risk:}
			\begin{itemize}
				\item Influenced by interdependencies
				\item Shared assumptions and methodologies
				\item Models that can affect multiple outcomes
			\end{itemize}
			\item \textbf{Duplication:}
			\begin{itemize}
				\item Similar models used in different parts
				\item Different assumptions
				\item Inefficiency and risk
			\end{itemize}
		\end{enumerate}
		
		\textcolor{myblue}{\textbf{Key Insight:}} Model landscape shows \textcolor{myblue}{how models interact and affect each other}!
	\end{frame}
	
	\begin{frame}{Slide 23: AI/ML Model Registration Considerations}
		\fontsize{6.5}{7.5}\selectfont
		\textbf{Registering AI/ML Models - New Aspects:}
		
		\vspace{0.1cm}
		\textbf{1. Complexity of Methodology and Design:}
		\begin{itemize}
			\item AI/ML models learn autonomously in developer-defined range
			\item Consider: chosen methodologies (neural networks vs. linear regression)
			\item Interpretability indicators (number of hidden layers, parameters)
		\end{itemize}
		
		\vspace{0.1cm}
		\textbf{2. Data Usage:}
		\begin{itemize}
			\item Volume of required data or number of features
			\item Complexity of data structures (unlabeled, metadata)
			\item Data quality (poorly labeled, low quality, unstructured)
			\item Variable interactions and transformations
		\end{itemize}
		
		\vspace{0.1cm}
		\textbf{3. Output Parameters:}
		\begin{itemize}
			\item Supervised vs. unsupervised learning
			\item Prediction vs. sentiment analysis vs. recommendations
			\item No direct way to evaluate output accuracy for unsupervised
		\end{itemize}
		
		\textcolor{myblue}{\textbf{Key Insight:}} AI/ML registration requires \textcolor{myblue}{additional considerations beyond traditional models}!
	\end{frame}
	
	\begin{frame}{Slide 24: Model Risk Ranking}
		\fontsize{6.5}{7.5}\selectfont
		\textbf{Model Risk Ranking Factors:}
		
		\vspace{0.1cm}
		\textbf{Key Factors:}
		\begin{enumerate}
			\item \textbf{Materiality:}
			\begin{itemize}
				\item Economic consequences of misuse
				\item Operational impact
				\item Reputational consequences
				\item Usage by different parties
				\item Impact on decisions, financial statements, or regulatory reporting
			\end{itemize}
			\item \textbf{Complexity:}
			\begin{itemize}
				\item May elevate risk ranking
				\item Even if materiality is low
				\item AI/ML models often have high complexity
			\end{itemize}
			\item \textbf{Exposure:}
			\begin{itemize}
				\item If model affects published financial reports
				\item Risk ranking may be elevated
				\item Even simple models can have high exposure
			\end{itemize}
		\end{enumerate}
		
		\vspace{0.1cm}
		\textbf{Important Note:}
		\begin{itemize}
			\item Minimum standards of model risk assessment should apply to all models
			\item Regardless of materiality
			\item Absence of assessment introduces its own model risk
		\end{itemize}
		
		\textcolor{myblue}{\textbf{Key Insight:}} Risk ranking considers \textcolor{myblue}{materiality, complexity, and exposure}!
	\end{frame}
	
	\begin{frame}{Slide 25: Model Performance Monitoring Metrics}
		\fontsize{6.5}{7.5}\selectfont
		\textbf{Performance Monitoring Metrics:}
		
		\vspace{0.1cm}
		\textbf{Classification Models:}
		\begin{itemize}
			\item Accuracy
			\item Precision
			\item Recall
			\item F1 Score
			\item Area Under the ROC Curve (AUC)
			\item Confusion Matrices
		\end{itemize}
		
		\vspace{0.1cm}
		\textbf{Regression Models:}
		\begin{itemize}
			\item Mean Absolute Error (MAE)
			\item Mean Squared Error (MSE)
			\item Coefficient of Determination (R-squared)
			\item Root Mean Square Error (RMSE)
		\end{itemize}
		
		\vspace{0.1cm}
		\textbf{Additional Metrics:}
		\begin{itemize}
			\item Clustering metrics
			\item Deep learning metrics
			\item GenAI performance metrics (see Chapter 10, Module 2)
		\end{itemize}
		
		\vspace{0.1cm}
		\textbf{Key Point:}
		\begin{itemize}
			\item Performance monitoring must be tailored to model type
			\item Metric must be appropriate for the model
			\item Model validation should assess this
		\end{itemize}
		
		\textcolor{myblue}{\textbf{Key Insight:}} Performance metrics must be \textcolor{myblue}{appropriate for the specific model type}!
	\end{frame}
	
	\begin{frame}{Slide 26: Three Lines of Defense - Overview}
		\fontsize{6.5}{7.5}\selectfont
		\textbf{Three Lines of Defense in Model Risk Management:}
		
		\vspace{0.1cm}
		\textbf{First Line (Business and Model Owners/Developers):}
		\begin{itemize}
			\item Owns the model and its performance
			\item Responsible for designing and implementing models
			\item Operates daily/weekly performance dashboards
			\item Owns data quality for model inputs
			\item First to react to threshold breaches
		\end{itemize}
		
		\vspace{0.1cm}
		\textbf{Second Line (Independent Risk Management/Model Validation):}
		\begin{itemize}
			\item Sets model risk management policy
			\item Conducts independent model validations
			\item Challenges First Line's performance assessments
			\item Maintains enterprise-wide model inventory
		\end{itemize}
		
		\vspace{0.1cm}
		\textbf{Third Line (Internal Audit):}
		\begin{itemize}
			\item Provides independent assurance to the board
			\item Reviews whether monitoring is appropriately archived
			\item Ensures policies are consistently followed
		\end{itemize}
		
		\textcolor{myblue}{\textbf{Key Insight:}} Three lines provide \textcolor{myblue}{checks and balances in model risk management}!
	\end{frame}
	
	\begin{frame}{Slide 27: First Line - Model Owners and Developers}
		\fontsize{6.5}{7.5}\selectfont
		\textbf{First Line of Defense - Detailed Responsibilities:}
		
		\vspace{0.1cm}
		\textbf{Model Owners:}
		\begin{itemize}
			\item Accountable for proper development, validation, approval
			\item Findings remediation, implementation, use
			\item Ongoing performance monitoring and reporting
			\item Change management
			\item Capture model characterization in inventory
			\item Assess model materiality
		\end{itemize}
		
		\vspace{0.1cm}
		\textbf{Model Developers:}
		\begin{itemize}
			\item Central to creation of models
			\item Decisions throughout modeling cycle affect risk
			\item Proper documentation saves time
			\item Act as both producers and recipients of validation feedback
		\end{itemize}
		
		\vspace{0.1cm}
		\textbf{AI/ML Considerations:}
		\begin{itemize}
			\item Ensure continuous learning models don't drift into unvalidated states
			\item Work closely with MLOps teams
			\item Data scientists must interpret metrics in business context
		\end{itemize}
		
		\textcolor{myblue}{\textbf{Key Insight:}} Model owners and developers are \textcolor{myblue}{responsible for model performance and accountability}!
	\end{frame}
	
	\begin{frame}{Slide 28: Second Line - Model Risk Management}
		\fontsize{6.5}{7.5}\selectfont
		\textbf{Second Line of Defense - Model Risk Function:}
		
		\vspace{0.1cm}
		\textbf{Key Responsibilities:}
		\begin{enumerate}
			\item \textbf{Framework Development:}
			\begin{itemize}
				\item Develop and maintain model risk management framework
				\item Establish general validation standards
			\end{itemize}
			\item \textbf{Model Inventory:}
			\begin{itemize}
				\item Maintain model inventory
				\item Combined with material aspects of model landscape
			\end{itemize}
			\item \textbf{Materiality Thresholds:}
			\begin{itemize}
				\item Set and maintain thresholds for model materiality
				\item Related to business decisions
			\end{itemize}
			\item \textbf{Validation Activities:}
			\begin{itemize}
				\item Ensure risk model validation activities occur
				\item Report validation results to board
			\end{itemize}
		\end{enumerate}
		
		\vspace{0.1cm}
		\textbf{AI/ML Considerations:}
		\begin{itemize}
			\item Deep expertise in validating complex ML models
			\item Assessing novel performance metrics
			\item Data drift detection for high-dimensional data
			\item Fairness/bias assessments
			\item Tools to "look inside the black box"
		\end{itemize}
		
		\textcolor{myblue}{\textbf{Key Insight:}} Second Line provides \textcolor{myblue}{independent challenge and oversight} of models!
	\end{frame}
	
	\begin{frame}{Slide 29: Third Line - Internal Audit}
		\fontsize{6.5}{7.5}\selectfont
		\textbf{Third Line of Defense - Internal Audit:}
		
		\vspace{0.1cm}
		\textbf{Key Responsibilities:}
		\begin{enumerate}
			\item \textbf{Independent Assurance:}
			\begin{itemize}
				\item Provide assurance to the board
				\item Model risk management framework is well-designed
				\item Operating effectively
			\end{itemize}
			\item \textbf{Annual Reviews:}
			\begin{itemize}
				\item Monitoring evidence appropriately archived
				\item Performance exceptions escalated according to policy
				\item Policies consistently followed and up-to-date
			\end{itemize}
		\end{enumerate}
		
		\vspace{0.1cm}
		\textbf{AI/ML Considerations:}
		\begin{itemize}
			\item Audit teams need sufficient understanding of ML concepts
			\item Focus on robustness of controls over automated ML pipelines
			\item Data governance for ML training data
			\item Adequacy of Second Line challenge for novel ML approaches
		\end{itemize}
		
		\vspace{0.1cm}
		\textbf{Key Point:}
		\begin{itemize}
			\item Audit is "meta-validation" or "validation of validation"
			\item Not the primary model risk management function
			\item Vital and necessary independent check
		\end{itemize}
		
		\textcolor{myblue}{\textbf{Key Insight:}} Internal Audit provides \textcolor{myblue}{independent verification} of the entire MRM framework!
	\end{frame}
	
	\begin{frame}{Slide 30: Communication and Interaction}
		\fontsize{6.5}{7.5}\selectfont
		\textbf{Communication and Interaction in Model Governance:}
		
		\vspace{0.1cm}
		\textbf{Key Areas:}
		\begin{enumerate}
			\item \textbf{Breach Notification:}
			\begin{itemize}
				\item How does First Line formally communicate performance breaches?
				\item What information must be provided?
			\end{itemize}
			\item \textbf{Independent Review:}
			\begin{itemize}
				\item Second Line's responsibilities and timelines
				\item Verifying issues and assessing materiality
				\item Challenging root cause analysis and remediation plans
			\end{itemize}
			\item \textbf{Dispute Resolution:}
			\begin{itemize}
				\item If First and Second Lines disagree
				\item Documented escalation path
				\item Model Risk Committee or Chief Risk Officer
			\end{itemize}
			\item \textbf{Reporting:}
			\begin{itemize}
				\item How findings flow to senior management and board?
				\item Consolidated and understandable manner
			\end{itemize}
		\end{enumerate}
		
		\vspace{0.1cm}
		\textbf{Common Failure Points:}
		\begin{itemize}
			\item First Line incentivized to downplay poor performance
			\item Second Line lacking business context or technical depth
			\item Third Line focusing on process rather than substance
		\end{itemize}
		
		\textcolor{myblue}{\textbf{Key Insight:}} Effective communication is \textcolor{myblue}{vital for successful model governance}!
	\end{frame}
	
	% ============================================================
	% SECTION 4: MODEL DEVELOPMENT AND TESTING - SLIDES 31-45
	% ============================================================
	
	\begin{frame}{Slide 31: Model Development Overview}
		\fontsize{6.5}{7.5}\selectfont
		\textbf{Model Development Process - Overview:}
		
		\vspace{0.1cm}
		\textbf{Key Steps:}
		\begin{enumerate}
			\item \textbf{Define Objectives and Scope:}
			\begin{itemize}
				\item Clear problem or risk being addressed
				\item Applicable products
				\item Required data and desired outcomes
			\end{itemize}
			\item \textbf{Data Collection and Preprocessing:}
			\begin{itemize}
				\item Gather relevant data
				\item Clean and transform data
				\item Handle missing values or outliers
			\end{itemize}
			\item \textbf{Data Analysis and Visualization:}
			\begin{itemize}
				\item Gain insights into data
				\item Identify patterns or relationships
				\item Understand limitations or biases
			\end{itemize}
			\item \textbf{Feature Engineering:}
			\begin{itemize}
				\item Select and engineer appropriate features
				\item Transform or combine variables
				\item Dimensionality reduction
			\end{itemize}
		\end{enumerate}
		
		\textcolor{myblue}{\textbf{Key Insight:}} Model development is a \textcolor{myblue}{systematic process from objectives to implementation}!
	\end{frame}
	
	\begin{frame}{Slide 32: Model Development Steps (Continued)}
		\fontsize{6.5}{7.5}\selectfont
		\textbf{Model Development Process (Continued):}
		
		\vspace{0.1cm}
		\textbf{5. Model Selection:}
		\begin{itemize}
			\item Choose modeling approach
			\item Comprehensive evaluation of alternatives
			\item Document rationale
		\end{itemize}
		
		\vspace{0.1cm}
		\textbf{6. Model Training/Calibration:}
		\begin{itemize}
			\item Optimize parameters and hyperparameters
			\item Best fit data and minimize errors
			\item Avoid data leakage (using future or concurrent data)
			\item Use cross-validation
		\end{itemize}
		
		\vspace{0.1cm}
		\textbf{7. Developer Testing:}
		\begin{itemize}
			\item Thorough testing before deployment
			\item Evaluate performance
			\item Identify issues early
		\end{itemize}
		
		\vspace{0.1cm}
		\textbf{8. Documentation:}
		\begin{itemize}
			\item Clear statement of purpose
			\item Sound design, theory, and logic
			\item Robust methodologies
			\item Rigorous data quality assessment
		\end{itemize}
		
		\textcolor{myblue}{\textbf{Key Insight:}} SR 11-7 requires \textcolor{myblue}{sound development process and thorough testing}!
	\end{frame}
	
	\begin{frame}{Slide 33: SR 11-7 Guidance on Model Development}
		\fontsize{6.5}{7.5}\selectfont
		\textbf{SR 11-7 Guidance on Model Development:}
		
		\vspace{0.1cm}
		\textbf{Key Requirements:}
		\begin{itemize}
			\item Clear statement of purpose to ensure model is developed in line with its intended use
			\item Sound design, theory, and logic underlying the model
			\item Robust model methodologies and processing components
			\item Rigorous assessment of data quality and relevance
			\item Appropriate documentation
		\end{itemize}
		
		\vspace{0.1cm}
		\textbf{Testing Requirements:}
		\begin{itemize}
			\item Integral part of model development
			\item Various components and overall functioning evaluated
			\item Show model is performing as intended
			\item Demonstrate accuracy, robustness, and stability
			\item Evaluate limitations and assumptions
		\end{itemize}
		
		\vspace{0.1cm}
		\textbf{Important Note:}
		\begin{itemize}
			\item Organizations should ensure development of judgmental and qualitative aspects is also sound
		\end{itemize}
		
		\textcolor{myblue}{\textbf{Key Insight:}} SR 11-7 requires \textcolor{myblue}{comprehensive development and testing of all model aspects}!
	\end{frame}
	
	\begin{frame}{Slide 34: Software Testing Types}
		\fontsize{6.5}{7.5}\selectfont
		\textbf{Types of Software Testing:}
		
		\vspace{0.1cm}
		\textbf{Testing Chronological Order:}
		\begin{enumerate}
			\item \textbf{Unit Tests:}
			\begin{itemize}
				\item Validate each piece of code performs as designed
			\end{itemize}
			\item \textbf{Component Tests:}
			\begin{itemize}
				\item Verify functionality of individual code sections
			\end{itemize}
			\item \textbf{Integration Tests:}
			\begin{itemize}
				\item Verify interfaces between components
			\end{itemize}
			\item \textbf{System Tests:}
			\begin{itemize}
				\item Ensure completely integrated system meets requirements
			\end{itemize}
			\item \textbf{Performance Tests:}
			\begin{itemize}
				\item Assess speed, responsiveness, and stability
			\end{itemize}
			\item \textbf{Regression Tests:}
			\begin{itemize}
				\item Verify system continues to function after modifications
			\end{itemize}
			\item \textbf{Acceptance Tests:}
			\begin{itemize}
				\item Determine if model is ready for deployment
			\end{itemize}
			\item \textbf{Adversarial Testing:}
			\begin{itemize}
				\item "Red team" evaluation of AI models
				\item Extreme scenarios and edge cases
				\item Prompt injection and unexpected queries
			\end{itemize}
		\end{enumerate}
		
		\textcolor{myblue}{\textbf{Key Insight:}} Testing should be \textcolor{myblue}{comprehensive and start as early as possible}!
	\end{frame}
	
	\begin{frame}{Slide 35: White-Box vs Black-Box Testing}
		\fontsize{6.5}{7.5}\selectfont
		\textbf{Testing Approaches:}
		
		\vspace{0.1cm}
		\textbf{White-Box Testing:}
		\begin{itemize}
			\item Testers have access to internal data structures, algorithms, and code
			\item Testing at unit level from developer's perspective
			\item Line-by-line proofreading of code
			\item More effective for finding specific issues
		\end{itemize}
		
		\vspace{0.1cm}
		\textbf{Black-Box Testing:}
		\begin{itemize}
			\item Treats software as a closed box
			\item From user's perspective
			\item No knowledge of internal implementation
			\item Reduces likelihood of bias
		\end{itemize}
		
		\vspace{0.1cm}
		\textbf{Gray-Box Testing:}
		\begin{itemize}
			\item Between white-box and black-box
			\item Informed by knowledge about functional specifications
			\item Knowledge of architecture of underlying code
			\item Does not access the code itself
		\end{itemize}
		
		\vspace{0.1cm}
		\textbf{When to Use:}
		\begin{itemize}
			\item White-box: internally developed models, early development
			\item Black-box: proprietary vendor models, functional testing
			\item Gray-box: distributed applications, testing component integration
		\end{itemize}
		
		\textcolor{myblue}{\textbf{Key Insight:}} Testing approach depends on \textcolor{myblue}{model type, development stage, and accessibility}!
	\end{frame}
	
	\begin{frame}{Slide 36: Test Configurations}
		\fontsize{6.5}{7.5}\selectfont
		\textbf{Effective Test Configurations:}
		
		\vspace{0.1cm}
		\textbf{Key Principles:}
		\begin{itemize}
			\item Tests need to be appropriately configured
			\item Usefulness hinges on adequacy and sufficiency of configurations
			\item Testers should not limit configurations solely to artificial cases
			\item Most valuable test configurations emerge from real-world situations
		\end{itemize}
		
		\vspace{0.1cm}
		\textbf{Less Realistic Parameters (Better Testing):}
		\begin{itemize}
			\item Missing trades
			\item Excessive trades
			\item Incorrect scales or mappings
			\item Extreme parameter values
			\item Parameters estimated from insufficient data
			\item Data provided by inexperienced users
			\item Variations in compilers or hardware
		\end{itemize}
		
		\vspace{0.1cm}
		\textbf{Recommendation:}
		\begin{itemize}
			\item Establish preliminary process that automatically generates test results from available data
		\end{itemize}
		
		\textcolor{myblue}{\textbf{Key Insight:}} The less realistic the test parameters, \textcolor{myblue}{the better the testing}!
	\end{frame}
	
	\begin{frame}{Slide 37: The Use Test}
		\fontsize{6.5}{7.5}\selectfont
		\textbf{The Use Test - Qualitative Assessment:}
		
		\vspace{0.1cm}
		\textbf{What is the Use Test?}
		\begin{itemize}
			\item Examines model within its context
			\item Considers human interaction and actual usage
			\item Assesses acceptance of the model
			\item Evaluates interpretation and application of results
		\end{itemize}
		
		\vspace{0.1cm}
		\textbf{Key Focus Areas:}
		\begin{enumerate}
			\item \textbf{People, Not Numbers:}
			\begin{itemize}
				\item Focuses on credible stories of acceptance
				\item Success and failure stories
				\item Technical challenges take a back seat
			\end{itemize}
			\item \textbf{Crucial Questions:}
			\begin{itemize}
				\item Are results utilized appropriately?
				\item Do they prove useful?
				\item If not, why?
			\end{itemize}
			\item \textbf{Exploration Areas:}
			\begin{itemize}
				\item Acceptance and trust
				\item Comprehensibility of results
				\item Communication of feedback to modelers
			\end{itemize}
		\end{enumerate}
		
		\textcolor{myblue}{\textbf{Key Insight:}} Use test focuses on \textcolor{myblue}{how people interact with and use the model}!
	\end{frame}
	
	\begin{frame}{Slide 38: Use Test Reporting}
		\fontsize{6.5}{7.5}\selectfont
		\textbf{Use Test - Reporting and Impact:}
		
		\vspace{0.1cm}
		\textbf{Reporting:}
		\begin{itemize}
			\item Results typically presented to senior management
			\item Not documented as technical reports or spreadsheets
			\item Easier to communicate results
			\item Minimizes risk of being seen as arduous requirement
		\end{itemize}
		
		\vspace{0.1cm}
		\textbf{Organizational Impact:}
		\begin{itemize}
			\item Contributes to improving culture of risk modeling
			\item Enhances risk management culture
			\item Bridges gap between technical and business perspectives
		\end{itemize}
		
		\vspace{0.1cm}
		\textbf{Benefits:}
		\begin{itemize}
			\item Promotes acceptance and trust
			\item Identifies usability issues
			\item Improves model communication
			\item Drives continuous improvement
		\end{itemize}
		
		\vspace{0.1cm}
		\textbf{Key Point:}
		\begin{itemize}
			\item Use test is qualitative in nature
			\item Difficult to follow a standard procedure
			\item Requires ongoing commitment
		\end{itemize}
		
		\textcolor{myblue}{\textbf{Key Insight:}} Use test reporting focuses on \textcolor{myblue}{senior management communication and culture improvement}!
	\end{frame}
	
	\begin{frame}{Slide 39: Limits on Model Use}
		\fontsize{6.5}{7.5}\selectfont
		\textbf{Should There Be Limits on Model Use?}
		
		\vspace{0.1cm}
		\textbf{Why Limits Are Needed:}
		\begin{itemize}
			\item Every QRM is based on underlying assumptions
			\item Assumptions should be clearly stated
			\item Identify model limitations and weaknesses
			\item Prevent misuse
		\end{itemize}
		
		\vspace{0.1cm}
		\textbf{Types of Limits:}
		\begin{enumerate}
			\item \textbf{Product Applicability:}
			\begin{itemize}
				\item Model approved for specific products only
				\item Example: Standard interest rate swaptions only
				\item Not for Bermudan swaption without additional validation
			\end{itemize}
			\item \textbf{Access Restrictions:}
			\begin{itemize}
				\item Who can use the model
				\item Training requirements
				\item Authorization levels
			\end{itemize}
			\item \textbf{Portfolio Size Limits:}
			\begin{itemize}
				\item Particularly for new markets or instruments
				\item Limited data or experience
				\item Risk ranking can change as portfolio grows
			\end{itemize}
		\end{enumerate}
		
		\textcolor{myblue}{\textbf{Key Insight:}} Model limits prevent \textcolor{myblue}{misuse and ensure risk ranking stays appropriate}!
	\end{frame}
	
	\begin{frame}{Slide 40: Model Use Guidelines}
		\fontsize{6.5}{7.5}\selectfont
		\textbf{Three General Guidelines for Risk Modeling Culture:}
		
		\vspace{0.1cm}
		\textbf{1. Awareness:}
		\begin{itemize}
			\item Be aware of limitations and assumptions of risk modeling
			\item Understand your company's history with QRMs
			\item Know risks entailed and validation processes in place
			\item Stay informed about market practices
			\item Recognize that the world is constantly changing
		\end{itemize}
		
		\vspace{0.1cm}
		\textbf{2. Transparency:}
		\begin{itemize}
			\item Transparently communicate assumptions, limitations, and documentation
			\item Provide detailed documentation with executive summaries
			\item Document decision-making process
			\item Document all validation activities, including unsuccessful attempts
			\item Engage in open communication with end users
		\end{itemize}
		
		\vspace{0.1cm}
		\textbf{3. Experience:}
		\begin{itemize}
			\item Learn from past modeling endeavors
			\item Apply relevant lessons
			\item Emphasize proper project management
			\item Develop prototypes early
			\item Establish and maintain libraries of reusable code
			\item Seek input from other modelers and external experts
		\end{itemize}
		
		\textcolor{myblue}{\textbf{Key Insight:}} Guidelines foster \textcolor{myblue}{culture that views validation as opportunity for insight}!
	\end{frame}
	
	\begin{frame}{Slide 41: SR 11-7 Adaptation}
		\fontsize{6.5}{7.5}\selectfont
		\textbf{Adapting SR 11-7 Guidelines:}
		
		\vspace{0.1cm}
		\textbf{Key Principle:}
		\begin{itemize}
			\item SR 11-7 established regulatory guidance
			\item Each institution needs to adapt to specific situation
		\end{itemize}
		
		\vspace{0.1cm}
		\textbf{Factors to Consider:}
		\begin{enumerate}
			\item \textbf{Specific Models:}
			\begin{itemize}
				\item Type of models used
				\item Complexity and risk profile
			\end{itemize}
			\item \textbf{Institution Size and Structure:}
			\begin{itemize}
				\item Large vs. small institution
				\item Organizational structure
				\item Resources available
			\end{itemize}
			\item \textbf{Governance Framework:}
			\begin{itemize}
				\item Existing governance structures
				\item Board and committee oversight
			\end{itemize}
			\item \textbf{Regulatory Requirements:}
			\begin{itemize}
				\item Specific regulatory expectations
				\item Jurisdictional requirements
			\end{itemize}
			\item \textbf{Budget Constraints:}
			\begin{itemize}
				\item Resources allocated to model risk
				\item Cost-benefit considerations
			\end{itemize}
		\end{enumerate}
		
		\textcolor{myblue}{\textbf{Key Insight:}} SR 11-7 is general guidance that \textcolor{myblue}{must be adapted to each institution's situation}!
	\end{frame}
	
	\begin{frame}{Slide 42: Model Changes Process}
		\fontsize{6.5}{7.5}\selectfont
		\textbf{Model Changes - Governance Process:}
		
		\vspace{0.1cm}
		\textbf{1. Change Management Process:}
		\begin{itemize}
			\item Structured process for making changes
			\item Steps: change requests, impact assessment, documentation, validation, approval
		\end{itemize}
		
		\vspace{0.1cm}
		\textbf{2. Documentation:}
		\begin{itemize}
			\item Comprehensive documentation required
			\item Reason for change, modifications, validation results, approvals
		\end{itemize}
		
		\vspace{0.1cm}
		\textbf{3. Independent Review:}
		\begin{itemize}
			\item Ongoing validation of changes
			\item Performed by individuals not involved in modification
			\item Assess accuracy, reliability, compliance, and risks
		\end{itemize}
		
		\vspace{0.1cm}
		\textbf{4. Approval:}
		\begin{itemize}
			\item Final approval through defined governance process
			\item Detailed documentation in model inventory
			\item Rationale, decision, conditions, or limitations
		\end{itemize}
		
		\vspace{0.1cm}
		\textbf{5. Review of Proposed Changes:}
		\begin{itemize}
			\item Model owner communicates change to MRM team
			\item Assessment for materiality
			\item Change-based validation if material
		\end{itemize}
		
		\textcolor{myblue}{\textbf{Key Insight:}} Model changes require \textcolor{myblue}{structured governance and independent review}!
	\end{frame}
	
	\begin{frame}{Slide 43: AI/ML Model Review Differences}
		\fontsize{6.5}{7.5}\selectfont
		\textbf{AI/ML Review Framework Differences:}
		
		\vspace{0.1cm}
		\textbf{1. Frequency:}
		\begin{itemize}
			\item More frequent monitoring for AI/ML models
			\item Higher complexity with more parameters
			\item Undesirable outcomes under certain conditions
			\item Greater susceptibility to data drift
			\item Monitoring dashboards may be useful
		\end{itemize}
		
		\vspace{0.1cm}
		\textbf{2. Validation Stakeholders:}
		\begin{itemize}
			\item Wider range of stakeholders affected
			\item Data protection regulation and ethical policies
			\item Include HR, compliance, operational risk
			\item Approval structure modifications needed
			\item Higher familiarity with AI/ML necessary
		\end{itemize}
		
		\vspace{0.1cm}
		\textbf{3. Validation Content:}
		\begin{itemize}
			\item Focus on explainability
			\item Benchmarking against traditional models
			\item Assessment based on different data sets
			\item Assessment of specific cases (edge cases)
			\item Back testing and cross-validation
			\item Reporting component for bias detection
		\end{itemize}
		
		\textcolor{myblue}{\textbf{Key Insight:}} AI/ML models require \textcolor{myblue}{more frequent review and broader stakeholder involvement}!
	\end{frame}
	
	\begin{frame}{Slide 44: Implementation and Monitoring}
		\fontsize{6.5}{7.5}\selectfont
		\textbf{Implementation and Monitoring of Model Changes:}
		
		\vspace{0.1cm}
		\textbf{Implementation:}
		\begin{itemize}
			\item Approved changes implemented in controlled manner
			\item Proper documentation
			\item Include changes to performance monitoring
		\end{itemize}
		
		\vspace{0.1cm}
		\textbf{Monitoring:}
		\begin{itemize}
			\item Ongoing monitoring continues after implementation
			\item Ensure changes achieve intended objectives
			\item Remain compliant with regulatory requirements
		\end{itemize}
		
		\vspace{0.1cm}
		\textbf{AI/ML Automation:}
		\begin{itemize}
			\item Growing volume of AI/ML models
			\item Potential need for more frequent retraining
			\item Greater impact of model errors
			\item Automated ongoing controls become more relevant
		\end{itemize}
		
		\vspace{0.1cm}
		\textbf{Key Point:}
		\begin{itemize}
			\item Automated controls where practical and responsible
			\item Balance automation with oversight
		\end{itemize}
		
		\textcolor{myblue}{\textbf{Key Insight:}} Model changes must be \textcolor{myblue}{implemented with controls and continuously monitored}!
	\end{frame}
	
	\begin{frame}{Slide 45: Model Changes Summary}
		\fontsize{6.5}{7.5}\selectfont
		\textbf{Model Changes - Key Takeaways:}
		
		\vspace{0.1cm}
		\textbf{Key Principles:}
		\begin{enumerate}
			\item \textbf{Structured Process:}
			\begin{itemize}
				\item Change requests, assessment, validation, approval
			\end{itemize}
			\item \textbf{Documentation:}
			\begin{itemize}
				\item Comprehensive record of all changes
				\item Rationale and approvals
			\end{itemize}
			\item \textbf{Independent Review:}
			\begin{itemize}
				\item Ongoing validation
				\item Objective assessment
			\end{itemize}
			\item \textbf{Frequency:}
			\begin{itemize}
				\item AI/ML models need more frequent review
				\item Address data drift and complexity
			\end{itemize}
			\item \textbf{Stakeholders:}
			\begin{itemize}
				\item Broader range of stakeholders
				\item Include compliance, HR, operational risk
			\end{itemize}
			\item \textbf{Monitoring:}
			\begin{itemize}
				\item Continuous monitoring after changes
				\item Automated controls where practical
			\end{itemize}
		\end{enumerate}
		
		\textcolor{myblue}{\textbf{Key Insight:}} Model changes require \textcolor{myblue}{comprehensive governance throughout the lifecycle}!
	\end{frame}
	
	% ============================================================
	% SECTION 5: MODEL VALIDATION - SLIDES 46-65
	% ============================================================
	
	\begin{frame}{Slide 46: Model Validation Definition}
		\fontsize{6.5}{7.5}\selectfont
		\textbf{What is Model Validation?}
		
		\vspace{0.1cm}
		\textbf{Definition:}
		\begin{itemize}
			\item Process of rigorously testing the model, inputs, outputs, governance, etc.
			\item To try to identify the extent of residual model risk
			\item Assess mitigating controls
		\end{itemize}
		
		\vspace{0.1cm}
		\textbf{Key Aspects:}
		\begin{itemize}
			\item Critical component in identifying model risk
			\item Involves qualitative and quantitative aspects
			\item Follows model lifecycle
		\end{itemize}
		
		\vspace{0.1cm}
		\textbf{When Validation Occurs:}
		\begin{enumerate}
			\item \textbf{Initial Validation:}
			\begin{itemize}
				\item Prior to implementation and usage
				\item Baseline assessment
			\end{itemize}
			\item \textbf{Routine Periodic Validation:}
			\begin{itemize}
				\item Annual reviews
				\item In-depth periodic baseline revalidations
			\end{itemize}
			\item \textbf{Change-Based Validation:}
			\begin{itemize}
				\item Triggered by material changes
				\item When model owner makes modifications
			\end{itemize}
		\end{enumerate}
		
		\textcolor{myblue}{\textbf{Key Insight:}} Model validation is the \textcolor{myblue}{rigorous testing to identify residual risk}!
	\end{frame}
	
	\begin{frame}{Slide 47: Validation Lifecycle}
		\fontsize{6.5}{7.5}\selectfont
		\textbf{Model Validation Lifecycle:}
		
		\vspace{0.1cm}
		\textbf{Key Stages:}
		\begin{enumerate}
			\item \textbf{Model Identification:}
			\begin{itemize}
				\item Recognize and classify the model
			\end{itemize}
			\item \textbf{Inventory:}
			\begin{itemize}
				\item Record in model inventory
				\item Schedule for validation
			\end{itemize}
			\item \textbf{Initial Validation:}
			\begin{itemize}
				\item Baseline validation prior to implementation
			\end{itemize}
			\item \textbf{Production and Monitoring:}
			\begin{itemize}
				\item Back testing and performance monitoring
				\item Ongoing throughout production usage
			\end{itemize}
			\item \textbf{Periodic Validation:}
			\begin{itemize}
				\item Routine validations
				\item Annual or more frequent
			\end{itemize}
			\item \textbf{Change-Based Validation:}
			\begin{itemize}
				\item Triggered by material changes
			\end{itemize}
			\item \textbf{Retirement:}
			\begin{itemize}
				\item Store in retired model inventory
				\item Decommission
			\end{itemize}
		\end{enumerate}
		
		\textcolor{myblue}{\textbf{Key Insight:}} Validation follows the \textcolor{myblue}{entire model lifecycle from identification to retirement}!
	\end{frame}
	
	\begin{frame}{Slide 48: Validation Activities}
		\fontsize{6.5}{7.5}\selectfont
		\textbf{Model Validation Activities Include:}
		
		\vspace{0.1cm}
		\textbf{Key Review Areas:}
		\begin{enumerate}
			\item \textbf{Data Assessment:}
			\begin{itemize}
				\item Appropriateness of data sources
				\item Data analysis (proxies, time series lengths)
			\end{itemize}
			\item \textbf{Assumptions and Parameters:}
			\begin{itemize}
				\item Assessment of any model assumptions
				\item Parameter evaluation
			\end{itemize}
			\item \textbf{Limitations:}
			\begin{itemize}
				\item Identification of limitations, strengths, weaknesses
				\item Approved use cases and users
			\end{itemize}
			\item \textbf{Implementation:}
			\begin{itemize}
				\item Appropriateness of implementation platform
				\item Usage of expert judgment and overlays
			\end{itemize}
			\item \textbf{Conceptual Soundness:}
			\begin{itemize}
				\item Review of conceptual soundness
				\item Assessment of benchmarks and best practices
			\end{itemize}
			\item \textbf{Performance Monitoring:}
			\begin{itemize}
				\item Analysis of ongoing monitoring plans
				\item Frequency, key risk indicators, escalation plans
			\end{itemize}
		\end{enumerate}
		
		\textcolor{myblue}{\textbf{Key Insight:}} Validation covers \textcolor{myblue}{all aspects of the model and its use}!
	\end{frame}
	
	\begin{frame}{Slide 49: Validation Activities (Continued)}
		\fontsize{6.5}{7.5}\selectfont
		\textbf{Model Validation Activities (Continued):}
		
		\vspace{0.1cm}
		\textbf{Additional Review Areas:}
		\begin{enumerate}
			\item \textbf{Documentation:}
			\begin{itemize}
				\item Review quality of documentation
				\item Ensure backup plans exist (internal and external)
			\end{itemize}
			\item \textbf{Controls:}
			\begin{itemize}
				\item Review controls around implementation and usage
				\item Access controls
			\end{itemize}
			\item \textbf{Outcomes Analysis:}
			\begin{itemize}
				\item Appropriateness of dependent variable
				\item Support for business objectives
			\end{itemize}
			\item \textbf{Testing:}
			\begin{itemize}
				\item Stress and scenario tests
				\item Model boundary testing
				\item Parameter stability
			\end{itemize}
			\item \textbf{Backtesting:}
			\begin{itemize}
				\item Appropriate cadence with model usage
				\item VaR models: daily monitoring
				\item Credit risk models: monthly
				\item Asset-liability models: quarterly
			\end{itemize}
		\end{enumerate}
		
		\textcolor{myblue}{\textbf{Key Insight:}} Validation includes \textcolor{myblue}{ongoing testing and monitoring appropriate to model type}!
	\end{frame}
	
	\begin{frame}{Slide 50: Initial Validation}
		\fontsize{6.5}{7.5}\selectfont
		\textbf{Initial Validation - Three Key Goals:}
		
		\vspace{0.1cm}
		\textbf{Goal 1: Operational Feasibility:}
		\begin{itemize}
			\item Checked and confirmed
			\item Model can run as intended without technical malfunctions
			\item Preproduction phase to run on intended schedule
		\end{itemize}
		
		\vspace{0.1cm}
		\textbf{Goal 2: Proper Documentation:}
		\begin{itemize}
			\item Adheres to firm-wide standards
			\item Model development standards
			\item Documentation standards
			\item Implementation standards
			\item Model monitoring standards
			\item Third-party governance standards
			\item Includes executive summaries
		\end{itemize}
		
		\vspace{0.1cm}
		\textbf{Goal 3: User Training:}
		\begin{itemize}
			\item Model users receive sufficient training
			\item Interpret and utilize results effectively
		\end{itemize}
		
		\vspace{0.1cm}
		\textbf{Key Point:}
		\begin{itemize}
			\item Results form basis for official approval (sign-off)
			\item Assumptions must be justified
			\item Limitations and weaknesses documented
		\end{itemize}
		
		\textcolor{myblue}{\textbf{Key Insight:}} Initial validation ensures \textcolor{myblue}{operational feasibility, documentation, and user training}!
	\end{frame}
	
	\begin{frame}{Slide 51: Ongoing Validation}
		\fontsize{6.5}{7.5}\selectfont
		\textbf{Ongoing (Periodic) Validation:}
		
		\vspace{0.1cm}
		\textbf{Primary Objective:}
		\begin{itemize}
			\item Observe whether model remains aligned with intended purpose
			\item Assumptions remain valid
			\item Data still appropriate
			\item Performance monitoring indicates expected performance
		\end{itemize}
		
		\vspace{0.1cm}
		\textbf{Key Activities:}
		\begin{enumerate}
			\item \textbf{Methodology Reassessment:}
			\begin{itemize}
				\item Still in line with best practices
				\item Reflects real world
			\end{itemize}
			\item \textbf{Assumption Validation:}
			\begin{itemize}
				\item Original assumptions remain valid
				\item Real world consistently challenges assumptions
			\end{itemize}
			\item \textbf{Benchmark/Challenger Models:}
			\begin{itemize}
				\item Support validation
				\item Compare results
			\end{itemize}
		\end{enumerate}
		
		\vspace{0.1cm}
		\textbf{Important Considerations:}
		\begin{itemize}
			\item Well documented
			\item Consolidated within single document
			\item Regular validation with fixed agenda
			\item Depth and frequency commensurate with risk tier
		\end{itemize}
		
		\textcolor{myblue}{\textbf{Key Insight:}} Ongoing validation ensures \textcolor{myblue}{models remain fit for purpose over time}!
	\end{frame}
	
	\begin{frame}{Slide 52: General Validation Considerations}
		\fontsize{6.5}{7.5}\selectfont
		\textbf{General Issues in Model Validation:}
		
		\vspace{0.1cm}
		\textbf{1. Terminology:}
		\begin{itemize}
			\item Alternative terms: model review, model QA, model evaluation, model vetting
		\end{itemize}
		
		\vspace{0.1cm}
		\textbf{2. Focus on Suitability:}
		\begin{itemize}
			\item Ultimate goal: assess if model is appropriate and effectively used
			\item Initial validation critical for credibility
			\item Ongoing performance monitoring as proposed and approved
			\item Validation must remain within intended application domain
		\end{itemize}
		
		\vspace{0.1cm}
		\textbf{3. Recursive Validation:}
		\begin{itemize}
			\item Validation activities should be subject to critical examination
			\item Ensure validation doesn't rely on defective data
			\item Must remain within intended application domain
			\item Need for challenge of validation activities
		\end{itemize}
		
		\vspace{0.1cm}
		\textbf{4. No "Valid" Model:}
		\begin{itemize}
			\item Models are simplified representations under assumed conditions
			\item Goal is not to test for a "valid" model
			\item Subject model to attempts to invalidate it
			\item Successful validation = withstood rigorous testing
		\end{itemize}
		
		\textcolor{myblue}{\textbf{Key Insight:}} Validation is about \textcolor{myblue}{testing, not proving validity}!
	\end{frame}
	
	\begin{frame}{Slide 53: Usefulness vs Validity}
		\fontsize{6.5}{7.5}\selectfont
		\textbf{Usefulness vs. Validity:}
		
		\vspace{0.1cm}
		\textbf{Key Principle:}
		\begin{itemize}
			\item Models do not have to be perfect to be useful
			\item If weaknesses and limitations are identified:
			\begin{itemize}
				\item Insights can improve the model
				\item Restrict application
				\item Enhance understanding of capabilities and limitations
			\end{itemize}
		\end{itemize}
		
		\vspace{0.1cm}
		\textbf{Important Considerations:}
		\begin{itemize}
			\item Model weaknesses can arise from regime shifts
			\item Need for ongoing monitoring and validation over time
		\end{itemize}
		
		\vspace{0.1cm}
		\textbf{Effective Challenge (SR 11-7):}
		\begin{itemize}
			\item Critical analysis by objective and informed parties
			\item Identify crucial model assumptions and limitations
			\item Spot and communicate relevant model weaknesses
		\end{itemize}
		
		\vspace{0.1cm}
		\textbf{Requirements for Effective Challenge:}
		\begin{itemize}
			\item Independence from model development
			\item High level of competence
			\item Sufficient degree of influence
		\end{itemize}
		
		\textcolor{myblue}{\textbf{Key Insight:}} Models don't need to be perfect to be \textcolor{myblue}{useful, but limitations must be understood}!
	\end{frame}
	
	\begin{frame}{Slide 54: GenAI Validation Considerations}
		\fontsize{6.5}{7.5}\selectfont
		\textbf{Validation of GenAI Models - Special Considerations:}
		
		\vspace{0.1cm}
		\textbf{Key Challenges:}
		\begin{enumerate}
			\item \textbf{Open-Ended Output:}
			\begin{itemize}
				\item Unstructured input and output
				\item Makes validation challenging
			\end{itemize}
			\item \textbf{Performance Metrics:}
			\begin{itemize}
				\item Different and more complex than traditional AI/ML
				\item New metrics needed
			\end{itemize}
			\item \textbf{Data Drift and Coverage:}
			\begin{itemize}
				\item Responses may become obsolete
				\item Changing environment
				\item Model obsolescence
			\end{itemize}
			\item \textbf{Hallucination:}
			\begin{itemize}
				\item Tendency to generate factually incorrect responses
				\item Appear very convincing
			\end{itemize}
			\item \textbf{Discrimination and Bias:}
			\begin{itemize}
				\item Ensure responses are not discriminatory or biased
				\item AI can amplify bias in data
			\end{itemize}
		\end{enumerate}
		
		\textcolor{myblue}{\textbf{Key Insight:}} GenAI validation requires \textcolor{myblue}{addressing unique challenges of generative models}!
	\end{frame}
	
	\begin{frame}{Slide 55: GenAI Validation Challenges (Continued)}
		\fontsize{6.5}{7.5}\selectfont
		\textbf{GenAI Validation Challenges (Continued):}
		
		\vspace{0.1cm}
		\textbf{6. Multi-Modal Issues:}
		\begin{itemize}
			\item Processing images, video, and audio data
			\item In addition to text
			\item More challenging to validate
		\end{itemize}
		
		\vspace{0.1cm}
		\textbf{7. Interpretability and Explainability:}
		\begin{itemize}
			\item Results difficult to explain or interpret
			\item Black-box nature
			\item Regulatory compliance challenges
		\end{itemize}
		
		\vspace{0.1cm}
		\textbf{8. Rapid Technological Change:}
		\begin{itemize}
			\item More frequent review required
			\item Regular recalibration needed
			\item Keeping pace with developments
		\end{itemize}
		
		\vspace{0.1cm}
		\textbf{9. Third-Party and Open-Source Models:}
		\begin{itemize}
			\item Increasingly offered as models as a service
			\item Particular attention needed
			\item Limited transparency
		\end{itemize}
		
		\vspace{0.1cm}
		\textbf{10. Computational Resources:}
		\begin{itemize}
			\item Significant resources required
			\item Training, validation, recalibration
			\item Infrastructure challenges
		\end{itemize}
		
		\textcolor{myblue}{\textbf{Key Insight:}} Human review is \textcolor{myblue}{essential in GenAI validation}!
	\end{frame}
	
	\begin{frame}{Slide 56: Model Design Motivation}
		\fontsize{6.5}{7.5}\selectfont
		\textbf{Model Design - Motivation for Development:}
		
		\vspace{0.1cm}
		\textbf{Changing Conditions:}
		\begin{itemize}
			\item New markets
			\item New products
			\item Customer requests
			\item Models must be validated
		\end{itemize}
		
		\vspace{0.1cm}
		\textbf{Regulatory Pressure:}
		\begin{itemize}
			\item Regulatory findings
			\item Model not performing as intended
			\item Redevelopment or replacement needed
		\end{itemize}
		
		\vspace{0.1cm}
		\textbf{Recent Crisis:}
		\begin{itemize}
			\item Losses highlight urgency
			\item Enhanced or more sensitive modeling needed
			\item Lessons from crisis
		\end{itemize}
		
		\vspace{0.1cm}
		\textbf{Internal Concerns:}
		\begin{itemize}
			\item Senior management discomfort
			\item Changes in business model
			\item Increased exposure to innovative products
			\item Changing market conditions
		\end{itemize}
		
		\vspace{0.1cm}
		\textbf{Innovation and Growth:}
		\begin{itemize}
			\item New or evolving customer demands
			\item New products or model extensions
			\item Business acquisitions
		\end{itemize}
		
		\textcolor{myblue}{\textbf{Key Insight:}} Model development is driven by \textcolor{myblue}{business needs, regulation, and changing conditions}!
	\end{frame}
	
	\begin{frame}{Slide 57: Model Developers}
		\fontsize{6.5}{7.5}\selectfont
		\textbf{Model Developers - The "Quants":}
		
		\vspace{0.1cm}
		\textbf{Required Expertise:}
		\begin{itemize}
			\item Mathematics
			\item Econometrics
			\item Statistics
			\item Strong product knowledge
		\end{itemize}
		
		\vspace{0.1cm}
		\textbf{Background:}
		\begin{itemize}
			\item Advanced academic degrees
			\item Mathematics, physics, engineering, computer science
			\item Scientific background
		\end{itemize}
		
		\vspace{0.1cm}
		\textbf{Key Qualities:}
		\begin{itemize}
			\item Creative problem-solving
			\item Skepticism
			\item Out-of-the-box thinking
			\item Experience including past crises
		\end{itemize}
		
		\vspace{0.1cm}
		\textbf{Important Balance:}
		\begin{itemize}
			\item Most challenging aspect: designing and constructing potential future scenarios
			\item Assigning weights to scenarios
			\item Must be mindful of difficulties in risk definition
			\item Avoid uncritical use of distributions
			\item Critical reflection on assumptions
		\end{itemize}
		
		\textcolor{myblue}{\textbf{Key Insight:}} Model developers need \textcolor{myblue}{quantitative expertise and practical experience}!
	\end{frame}
	
	\begin{frame}{Slide 58: Modeling Issues - Discretization}
		\fontsize{6.5}{7.5}\selectfont
		\textbf{Modeling Issues - Discretization:}
		
		\vspace{0.1cm}
		\textbf{What is Discretization?}
		\begin{itemize}
			\item Describing the world using continuous variables
			\item Cases where variables are discretized to finite values
			\item Reducing complexity in modeling path-dependent events
			\item Events occurring at discrete time intervals
		\end{itemize}
		
		\vspace{0.1cm}
		\textbf{Why Discretize?}
		\begin{itemize}
			\item Complex equations without closed-form solutions
			\item Numerical approximation via discretization
			\item Should yield very close results when properly implemented
			\item Time intervals divided into days, months, years
		\end{itemize}
		
		\vspace{0.1cm}
		\textbf{Challenges:}
		\begin{itemize}
			\item Discretization error
			\item Stability
			\item Convergence
			\item Computational complexity
			\item Boundary conditions
			\item Numerical precision
			\item Choice of time step is crucial
		\end{itemize}
		
		\textcolor{myblue}{\textbf{Key Insight:}} Discretization enables implementation but \textcolor{myblue}{requires careful parameter selection}!
	\end{frame}
	
	\begin{frame}{Slide 59: Modeling Issues - Approximations}
		\fontsize{6.5}{7.5}\selectfont
		\textbf{Modeling Issues - Approximations:}
		
		\vspace{0.1cm}
		\textbf{What are Approximations?}
		\begin{itemize}
			\item Replacing complex or hard-to-evaluate components
			\item Alternative methods that produce similar results
			\item More conveniently implemented
		\end{itemize}
		
		\vspace{0.1cm}
		\textbf{Key Considerations:}
		\begin{itemize}
			\item Choosing an approximation fixes the precision
			\item No room for improvement except by selecting different approximation
		\end{itemize}
		
		\vspace{0.1cm}
		\textbf{Examples:}
		\begin{itemize}
			\item Price sensitivities
			\item Taylor series approximation in market risk
			\item Valid for simple instruments and small changes
		\end{itemize}
		
		\vspace{0.1cm}
		\textbf{Time Horizon Scaling:}
		\begin{itemize}
			\item Scaling VaR by square root of time
			\item Assumes normal distribution and temporal independence
			\item Often convenient but rarely empirically justified
			\item Validity must be examined via statistical tests
		\end{itemize}
		
		\textcolor{myblue}{\textbf{Key Insight:}} Approximations simplify but \textcolor{myblue}{fix precision and require careful validation}!
	\end{frame}
	
	\begin{frame}{Slide 60: Modeling Issues - Numerical Evaluation}
		\fontsize{6.5}{7.5}\selectfont
		\textbf{Modeling Issues - Numerical Evaluation:}
		
		\vspace{0.1cm}
		\textbf{What is Numerical Evaluation?}
		\begin{itemize}
			\item Evaluation of model component with potential for arbitrary precision
			\item Equipped with parameters controlling precision
		\end{itemize}
		
		\vspace{0.1cm}
		\textbf{Methods:}
		\begin{itemize}
			\item Computation of integrals
			\item Evaluating non-elementary functions
			\item Solving systems of equations
			\item Linear algebra problems
			\item Interpolation
			\item Numerical differential equation solutions
			\item Optimization
		\end{itemize}
		
		\vspace{0.1cm}
		\textbf{Key Factors:}
		\begin{itemize}
			\item \textbf{Stability:} Small errors should not lead to unbounded imprecision
			\item \textbf{Convergence:} Results approach true value as precision increases
			\item \textbf{Floating-Point Arithmetic:} Rounding errors may affect stability
		\end{itemize}
		
		\vspace{0.1cm}
		\textbf{Trade-off:}
		\begin{itemize}
			\item Running time vs. precision
			\item Theoretical confirmation rarely sufficient
			\item Extensive testing necessary
		\end{itemize}
		
		\textcolor{myblue}{\textbf{Key Insight:}} Numerical evaluation requires \textcolor{myblue}{balancing precision with computational cost}!
	\end{frame}
	
	\begin{frame}{Slide 61: Modeling Issues - Random Numbers}
		\fontsize{6.5}{7.5}\selectfont
		\textbf{Modeling Issues - Random Numbers:}
		
		\vspace{0.1cm}
		\textbf{Why Random Numbers Needed:}
		\begin{itemize}
			\item Models rely heavily on probability theory
			\item Assigning values to random variables
			\item Samples drawn from distributions
		\end{itemize}
		
		\vspace{0.1cm}
		\textbf{Challenges:}
		\begin{itemize}
			\item Generating random numbers is difficult for both humans and computers
			\item Computers use pseudo-random number generators (RNGs)
			\item Depend on initial seed
			\item Cannot produce actual random numbers
		\end{itemize}
		
		\vspace{0.1cm}
		\textbf{Best Practice:}
		\begin{itemize}
			\item Rely on well-designed deterministic pseudo-RNGs
			\item Rather than introducing randomness in choice of RNG
		\end{itemize}
		
		\vspace{0.1cm}
		\textbf{RNG Quality Assessment:}
		\begin{itemize}
			\item Specialized test suites
			\item Criteria: speed and reproducibility
		\end{itemize}
		
		\textcolor{myblue}{\textbf{Key Insight:}} Random number generation requires \textcolor{myblue}{high-quality deterministic pseudo-RNGs}!
	\end{frame}
	
	\begin{frame}{Slide 62: Model Implementation Tasks}
		\fontsize{6.5}{7.5}\selectfont
		\textbf{Model Implementation Tasks:}
		
		\vspace{0.1cm}
		\textbf{1. Model Design:}
		\begin{itemize}
			\item Designed and documented by model designers
			\item Facilitate translation into computer code
		\end{itemize}
		
		\vspace{0.1cm}
		\textbf{2. Core Implementation:}
		\begin{itemize}
			\item Inner workings implemented as black box
			\item With specified interfaces
		\end{itemize}
		
		\vspace{0.1cm}
		\textbf{3. System Implementation:}
		\begin{itemize}
			\item Model core integrated into system
			\item Provides user interfaces
			\item Collects input data and parameters
			\item Schedules computations
			\item Feeds the core
			\item Processes output data
			\item Keeps history of computations
		\end{itemize}
		
		\vspace{0.1cm}
		\textbf{Key Considerations:}
		\begin{itemize}
			\item Different skills for core vs system implementation
			\item Different programming languages
			\item Specialized skills needed for system implementation
		\end{itemize}
		
		\textcolor{myblue}{\textbf{Key Insight:}} Implementation involves \textcolor{myblue}{model design, core development, and system integration}!
	\end{frame}
	
	\begin{frame}{Slide 63: Model Adaptation Tasks}
		\fontsize{6.5}{7.5}\selectfont
		\textbf{Model Adaptation Tasks:}
		
		\vspace{0.1cm}
		\textbf{Why Adaptation Needed:}
		\begin{itemize}
			\item No model lasts forever
			\item Periodic reviews detect weaknesses
			\item Demand adaptation or revision
		\end{itemize}
		
		\vspace{0.1cm}
		\textbf{Adaptation Types:}
		\begin{enumerate}
			\item \textbf{Model Adaptation:}
			\begin{itemize}
				\item Model and documentation adapted (or replaced)
				\item Existing computer code adapted
			\end{itemize}
			\item \textbf{Core Adaptation:}
			\begin{itemize}
				\item Model core adapted (or replaced)
				\item Interfaces may be adapted
			\end{itemize}
			\item \textbf{System Adaptation:}
			\begin{itemize}
				\item System around model core adapted (or replaced)
			\end{itemize}
		\end{enumerate}
		
		\vspace{0.1cm}
		\textbf{Challenges:}
		\begin{itemize}
			\item Initial implementation team may have left
			\item External consultants may have done some tasks
			\item If model no longer understood, adaptation is challenging
			\item If same people, risk of ignorance or entrenched poor practices
		\end{itemize}
		
		\textcolor{myblue}{\textbf{Key Insight:}} Adaptation requires \textcolor{myblue}{understanding the model and careful change management}!
	\end{frame}
	
	\begin{frame}{Slide 64: Running a Model}
		\fontsize{6.5}{7.5}\selectfont
		\textbf{Running a Model - Key Considerations:}
		
		\vspace{0.1cm}
		\textbf{Robustness:}
		\begin{itemize}
			\item Models can become complex
			\item Intertwining with multiple systems
			\item Increased risk of failure
			\item Minimize complexity
			\item Eliminate exotic systems
			\item Standardize interfaces
			\item Avoid "reinventing the wheel"
		\end{itemize}
		
		\vspace{0.1cm}
		\textbf{Computational Considerations:}
		\begin{itemize}
			\item Total response time varies greatly
			\item Credit risk models: may take weeks
			\item Market risk models: daily reporting
			\item Pre-deal checking: faster processing
			\item Overall running time determined by slowest task
			\item Over-optimization for specific hardware can be risky
		\end{itemize}
		
		\vspace{0.1cm}
		\textbf{Flexibility:}
		\begin{itemize}
			\item Scheduled runs: regular numbers with minimal manual intervention
			\item Ad hoc runs: additional information on short notice
			\item Staff knowledgeable about model interfaces
			\item Prevent gradual shift from scheduled to ad hoc mode
		\end{itemize}
		
		\textcolor{myblue}{\textbf{Key Insight:}} Running models requires \textcolor{myblue}{robustness, appropriate computing, and flexibility}!
	\end{frame}
	
	\begin{frame}{Slide 65: Model Adaptation and Running Summary}
		\fontsize{6.5}{7.5}\selectfont
		\textbf{Implementation, Adaptation, and Running - Key Takeaways:}
		
		\vspace{0.1cm}
		\textbf{Implementation:}
		\begin{itemize}
			\item Model design, core implementation, system integration
			\item Different skills for different tasks
			\item Documentation is essential
		\end{itemize}
		
		\vspace{0.1cm}
		\textbf{Adaptation:}
		\begin{itemize}
			\item Periodic reviews detect weaknesses
			\item Model, core, or system adaptation
			\item Challenges with changing teams
		\end{itemize}
		
		\vspace{0.1cm}
		\textbf{Running:}
		\begin{itemize}
			\item Robustness and standardization
			\item Computational considerations
			\item Scheduled vs. ad hoc runs
		\end{itemize}
		
		\vspace{0.1cm}
		\textbf{Common Issues:}
		\begin{itemize}
			\item Over-optimism
			\item Inadequate budgeting
			\item Midstream redefinition of goals
			\item Misunderstandings and miscommunication
			\item Programming errors
			\item Lack of software quality assurance
			\item Inappropriate design architecture
		\end{itemize}
		
		\textcolor{myblue}{\textbf{Key Insight:}} Model operations require \textcolor{myblue}{careful management of implementation, adaptation, and running}!
	\end{frame}
	
	% ============================================================
	% SECTION 6: MISINTERPRETATION OF RESULTS - SLIDES 66-75
	% ============================================================
	
	\begin{frame}{Slide 66: Misinterpretation of Model Results}
		\fontsize{6.5}{7.5}\selectfont
		\textbf{Potential Sources of Misinterpretation:}
		
		\vspace{0.1cm}
		\textbf{1. Lack of Comprehension:}
		\begin{itemize}
			\item Users lack comprehensive understanding of model
			\item Underlying assumptions unknown
			\item Limitations of outputs unclear
			\item Necessary knowledge and expertise missing
		\end{itemize}
		
		\vspace{0.1cm}
		\textbf{2. Neglecting Model Limitations:}
		\begin{itemize}
			\item Models have inherent limitations
			\item May assume certain conditions
			\item May not fully capture all relevant risks
			\item Failing to acknowledge limitations leads to flawed understanding
		\end{itemize}
		
		\vspace{0.1cm}
		\textbf{3. Overemphasis on Point Estimates:}
		\begin{itemize}
			\item Models provide estimates with uncertainty
			\item Focusing solely on point estimates
			\item Ignoring range of potential outcomes
			\item Understanding probability distributions is crucial
		\end{itemize}
		
		\textcolor{myblue}{\textbf{Key Insight:}} Misinterpretation can have \textcolor{myblue}{significant implications for decision-making}!
	\end{frame}
	
	\begin{frame}{Slide 67: Misinterpretation (Continued)}
		\fontsize{6.5}{7.5}\selectfont
		\textbf{Potential Sources of Misinterpretation (Continued):}
		
		\vspace{0.1cm}
		\textbf{4. Failure to Consider Qualitative Factors:}
		\begin{itemize}
			\item Models focus on quantitative measures
			\item Qualitative factors should not be overlooked
			\item Market conditions, regulatory changes, external influences
			\item Exclusion can lead to distorted understanding
		\end{itemize}
		
		\vspace{0.1cm}
		\textbf{5. Ignoring Decision Context:}
		\begin{itemize}
			\item Model results considered independently
			\item Without specific context of decision
			\item Models should inform decision making within context
			\item Incorporate broader business strategies, risk appetite
		\end{itemize}
		
		\vspace{0.1cm}
		\textbf{6. Confirmation Bias:}
		\begin{itemize}
			\item Users interpret results to align with preconceived beliefs
			\item Cloud judgment
			\item Compromise objective interpretation
			\item Approach results with open mind and commitment to unbiased analysis
		\end{itemize}
		
		\textcolor{myblue}{\textbf{Key Insight:}} Misinterpretation arises from \textcolor{myblue}{cognitive biases and lack of context}!
	\end{frame}
	
	\begin{frame}{Slide 68: Fostering Proper Interpretation}
		\fontsize{6.5}{7.5}\selectfont
		\textbf{Fostering Proper Interpretation of Model Results:}
		
		\vspace{0.1cm}
		\textbf{Key Practices:}
		\begin{enumerate}
			\item \textbf{Training and Education:}
			\begin{itemize}
				\item Ensure users understand model
				\item Explain assumptions and limitations
				\item Provide context for results
			\end{itemize}
			\item \textbf{Clear Communication:}
			\begin{itemize}
				\item Present results with uncertainty ranges
				\item Include confidence intervals
				\item Explain probabilistic interpretations
			\end{itemize}
			\item \textbf{Open Dialogue:}
			\begin{itemize}
				\item Encourage questioning of results
				\item Discuss alternative interpretations
				\item Foster critical thinking
			\end{itemize}
			\item \textbf{Culture:}
			\begin{itemize}
				\item Foster culture of open dialogue
				\item Encourage critical thinking
				\item Cross-functional collaboration
			\end{itemize}
		\end{enumerate}
		
		\vspace{0.1cm}
		\textbf{Key Point:}
		\begin{itemize}
			\item Guard against misinterpretation biases
			\item Approach results with open mind
			\item Commit to unbiased analysis
		\end{itemize}
		
		\textcolor{myblue}{\textbf{Key Insight:}} Proper interpretation requires \textcolor{myblue}{training, clear communication, and supportive culture}!
	\end{frame}
	







\begin{frame}{Slide 69: Model Governance for AI/ML - Recommendations}
	\fontsize{6.5}{7.5}\selectfont
	
	\textbf{Recommendations for AI/ML Model Governance:}
	
	\vspace{0.1cm}
	\textbf{1. Begin with Existing Frameworks:}
	\begin{itemize}
		\item Don't start from scratch
		\item Use experience from traditional model validation
		\item Build on a strong foundation
	\end{itemize}
	
	\vspace{0.1cm}
	\textbf{2. Consider the New Role of Data:}
	\begin{itemize}
		\item AI/ML is ``model-free'' -- everything depends on data
		\item Address issues such as:
		\begin{itemize}
			\item Various forms of bias
			\item Overfitting
			\item Population drift
			\item Model drift and regime changes
			\item Risk of harmful output
		\end{itemize}
		
		\item New uses will reveal further challenges:
		\begin{itemize}
			\item Data limitations
			\item Biases
			\item Privacy and right to use
			\item Legality
		\end{itemize}
	\end{itemize}
	
	\textcolor{myblue}{\textbf{Key Insight:}}
	AI/ML governance must
	\textcolor{myblue}{address the central role of data}!
\end{frame}








	
	\begin{frame}{Slide 70: AI/ML Governance Recommendations (Continued)}
		\fontsize{6.5}{7.5}\selectfont
		\textbf{Recommendations for AI/ML Model Governance (Continued):}
		
		\vspace{0.1cm}
		\textbf{3. Add New Perspectives:}
		\begin{itemize}
			\item AI/ML increases relevance of ethical aspects
			\item Data bias
			\item Explainability of model results
			\item Recalibration process
			\item These attributes need to be considered in model inventory
		\end{itemize}
		
		\vspace{0.1cm}
		\textbf{4. BIS Survey Findings - Data Governance:}
		\begin{itemize}
			\item Adequate documentation
			\item Policies, standards, procedures
			\item Data quality - accuracy, transparency, provenance
			\item Security and responsible use
			\item Comprehensive data management systems
			\item Metadata on sources, version, curation
			\item Asset inventories and metadata registries
		\end{itemize}
		
		\textcolor{myblue}{\textbf{Key Insight:}} AI/ML governance must \textcolor{myblue}{incorporate ethical considerations and robust data governance}!
	\end{frame}
	
	\begin{frame}{Slide 71: BIS Survey - Organizational Structures}
		\fontsize{6.5}{7.5}\selectfont
		\textbf{BIS Survey - Organizational Structures for AI Governance:}
		
		\vspace{0.1cm}
		\textbf{Common Approaches:}
		\begin{enumerate}
			\item \textbf{Centralized Solutions:}
			\begin{itemize}
				\item Steering committees/commissions
				\item Program management offices
				\item Multi-modal structures
			\end{itemize}
			\item \textbf{Hub-and-Spoke Systems:}
			\begin{itemize}
				\item Overall network with multiple functional hubs
				\item Connected to various departments
			\end{itemize}
			\item \textbf{Fully Decentralized:}
			\begin{itemize}
				\item Governance at business line level
				\item Local control
			\end{itemize}
		\end{enumerate}
		
		\vspace{0.1cm}
		\textbf{Accountability Assignments:}
		\begin{itemize}
			\item Existing roles: Chief Data/Information/Technology Officer
			\item Existing units: IT department
			\item Newly created functions
			\item Informal setups: communities of practice, cross-functional networks
		\end{itemize}
		
		\textcolor{myblue}{\textbf{Key Insight:}} AI governance structures range from \textcolor{myblue}{centralized to decentralized with various accountability models}!
	\end{frame}
	
	\begin{frame}{Slide 72: BIS Survey - Guidelines and Risk Management}
		\fontsize{6.5}{7.5}\selectfont
		\textbf{BIS Survey - Guidelines and Risk Management:}
		
		\vspace{0.1cm}
		\textbf{Guidelines and Policies:}
		\begin{itemize}
			\item Terms of use for AI systems
			\item Modalities, interfaces, and data
			\item Models, third-party dependencies
			\item Upstream/downstream data
			\item Guidelines for responsible AI use
			\item Privacy standards
			\item Explainability requirements
			\item International principles (e.g., Fundamental Principles of Official Statistics)
			\item Documentation of methods for regular assessment
		\end{itemize}
		
		\vspace{0.1cm}
		\textbf{Risk Management Frameworks:}
		\begin{itemize}
			\item Regular audit of AI systems for vulnerabilities
			\item Evaluation of robustness of safety measures
			\item Incident response procedures
			\item Business continuity plans
			\item Processes to promote and maintain skills
			\item Prevent erosion of knowledge due to overreliance
			\item Leverage existing frameworks (three lines of defense)
		\end{itemize}
		
		\textcolor{myblue}{\textbf{Key Insight:}} Comprehensive guidelines and risk management are \textcolor{myblue}{essential for AI governance}!
	\end{frame}
	
	\begin{frame}{Slide 73: Governance Challenges of GenAI}
		\fontsize{6.5}{7.5}\selectfont
		\textbf{Why GenAI Governance Is Harder:}
		
		\vspace{0.1cm}
		\textbf{1. Opacity and Unpredictability:}
		\begin{itemize}
			\item Functioning almost always opaque
			\item Output unpredictable
			\item Difficult to interpret or audit
			\item Probabilistic nature: same prompt may lead to different outputs
			\item Reproducibility challenges
		\end{itemize}
		
		\vspace{0.1cm}
		\textbf{2. Emergent Capabilities:}
		\begin{itemize}
			\item Behaviors not anticipated by developers
			\item As models become larger and more sophisticated
			\item Skills not observed in smaller models
			\item Changing capabilities hard to assess
		\end{itemize}
		
		\vspace{0.1cm}
		\textbf{3. Post-Deployment Shifts:}
		\begin{itemize}
			\item Model behavior can shift after deployment
			\item External providers may update or fine-tune models
			\item Changes in prompts or connected data alter performance
			\item One-off validation is inadequate
		\end{itemize}
		
		\textcolor{myblue}{\textbf{Key Insight:}} GenAI governance must address \textcolor{myblue}{opacity, unpredictability, and evolving capabilities}!
	\end{frame}
	
	\begin{frame}{Slide 74: GenAI Governance Approaches}
		\fontsize{6.5}{7.5}\selectfont
		\textbf{Internal Governance Levers for GenAI:}
		
		\vspace{0.1cm}
		\textbf{1. Assign Ownership and Accountability:}
		\begin{itemize}
			\item Clarify responsibility at each stage
			\item Tool selection, testing, deployment, monitoring
			\item Ensure responsibility doesn't fall between teams
		\end{itemize}
		
		\vspace{0.1cm}
		\textbf{2. Risk-Based Review and Deployment:}
		\begin{itemize}
			\item Tiered review gates
			\item Full risk reviews for high-impact systems
			\item Light checks for low-stakes internal tools
		\end{itemize}
		
		\vspace{0.1cm}
		\textbf{3. Incentivize Responsible Behavior:}
		\begin{itemize}
			\item Risk-tied performance metrics
			\item Reward flagging risks, documenting limitations
			\item Align team incentives with responsible use
		\end{itemize}
		
		\vspace{0.1cm}
		\textbf{4. Staff Training:}
		\begin{itemize}
			\item Equip staff with knowledge of GenAI risks
			\item Data protection, hallucinations, complaints
		\end{itemize}
		
		\textcolor{myblue}{\textbf{Key Insight:}} Internal governance requires \textcolor{myblue}{clear ownership, risk-based reviews, and proper incentives}!
	\end{frame}
	
	\begin{frame}{Slide 75: External Governance and Adaptive Tools}
		\fontsize{6.5}{7.5}\selectfont
		\textbf{External Governance and Adaptive Tools:}
		
		\vspace{0.1cm}
		\textbf{External Governance Constraints:}
		\begin{itemize}
			\item Track evolving legal and regulatory requirements
			\item EU AI Act, Digital Services Act, US state laws
			\item Evaluate responsibilities with model providers
			\item Understand incident policies and update schedules
			\item Align with emerging standards (NIST AI RMF, ISO/IEC 42001)
		\end{itemize}
		
		\vspace{0.1cm}
		\textbf{Adaptive Tools and Practices:}
		\begin{itemize}
			\item Continuously monitor post-deployment usage
			\item Use red-teaming and adversarial testing
			\item Develop scenario plans for unpredictability
			\item Deploy technical safeguards (watermarking, usage logging)
			\item Enable feedback loops and update mechanisms
		\end{itemize}
		
		\textcolor{myblue}{\textbf{Key Insight:}} GenAI governance requires \textcolor{myblue}{external awareness and adaptive practices}!
	\end{frame}
	
	% ============================================================
	% SECTION 7: QUESTIONS FOR RISK PROFESSIONALS - SLIDES 76-85
	% ============================================================
	
	\begin{frame}{Slide 76: Strategic Risk Assessment Questions}
		\fontsize{6.5}{7.5}\selectfont
		\textbf{Strategic Risk Assessment Questions:}
		
		\vspace{0.1cm}
		\textbf{Key Questions:}
		\begin{enumerate}
			\item \textbf{Is GenAI the right tool for this application?}
			\begin{itemize}
				\item Ensure reasons for deploying are good reasons
				\item Know what trade-offs are to be expected
			\end{itemize}
			\item \textbf{What are the risks to:}
			\begin{itemize}
				\item Safety
				\item Legality
				\item Reputation
				\item User trust
				\item Transparency
			\end{itemize}
			\item \textbf{Risk Identification:}
			\begin{itemize}
				\item Identify relevant risks before deployment
				\item Map potential unpredictable behaviors
				\item Assess possible impacts
			\end{itemize}
		\end{enumerate}
		
		\vspace{0.1cm}
		\textbf{Key Point:}
		\begin{itemize}
			\item Strategic assessment should occur before deployment
			\item Consider both intended and unintended consequences
		\end{itemize}
		
		\textcolor{myblue}{\textbf{Key Insight:}} Strategic assessment ensures \textcolor{myblue}{GenAI is the right tool for the right application}!
	\end{frame}
	
	\begin{frame}{Slide 77: Model Selection Questions}
		\fontsize{6.5}{7.5}\selectfont
		\textbf{Model Selection and Deployment Questions:}
		
		\vspace{0.1cm}
		\textbf{Key Questions:}
		\begin{enumerate}
			\item \textbf{What model is being used?}
			\begin{itemize}
				\item What do we know about its training data?
				\item What do we know about its alignment?
			\end{itemize}
			\item \textbf{Do we control the model?}
			\begin{itemize}
				\item API vs open source
				\item Control over updates
				\item Transparency varies
			\end{itemize}
		\end{enumerate}
		
		\vspace{0.1cm}
		\textbf{Why These Questions Matter:}
		\begin{itemize}
			\item Understanding training data reveals potential biases
			\item Control determines ability to audit and modify
			\item Transparency affects governance capabilities
		\end{itemize}
		
		\vspace{0.1cm}
		\textbf{Key Point:}
		\begin{itemize}
			\item Model selection affects all subsequent governance
			\item Choose models that align with governance capabilities
		\end{itemize}
		
		\textcolor{myblue}{\textbf{Key Insight:}} Model selection determines \textcolor{myblue}{governance options and risk exposure}!
	\end{frame}
	
	\begin{frame}{Slide 78: Data Governance Questions}
		\fontsize{6.5}{7.5}\selectfont
		\textbf{Data Governance Questions:}
		
		\vspace{0.1cm}
		\textbf{Key Questions:}
		\begin{enumerate}
			\item \textbf{What data is being used?}
			\begin{itemize}
				\item Is any of it proprietary or sensitive?
				\item Legal restrictions (GDPR, etc.)
				\item Client confidentiality
				\item Intellectual property
			\end{itemize}
			\item \textbf{Are data protection and retention policies in place?}
			\begin{itemize}
				\item Trace how inputs are used by providers
				\item Understand data flow
			\end{itemize}
		\end{enumerate}
		
		\vspace{0.1cm}
		\textbf{Why These Questions Matter:}
		\begin{itemize}
			\item Data provenance affects legal compliance
			\item Data protection policies prevent breaches
			\item Understanding data flow enables governance
		\end{itemize}
		
		\vspace{0.1cm}
		\textbf{Key Point:}
		\begin{itemize}
			\item Data governance is foundational for GenAI
			\item Know what data is used and how it's protected
		\end{itemize}
		
		\textcolor{myblue}{\textbf{Key Insight:}} Data governance questions are \textcolor{myblue}{essential for responsible GenAI use}!
	\end{frame}
	
	\begin{frame}{Slide 79: Technical Risk Control Questions}
		\fontsize{6.5}{7.5}\selectfont
		\textbf{Technical Risk Control Questions:}
		
		\vspace{0.1cm}
		\textbf{Key Questions:}
		\begin{enumerate}
			\item \textbf{Have we tested for unpredictability?}
			\begin{itemize}
				\item Hallucinations
				\item Model drift
				\item Unexpected behaviors
			\end{itemize}
			\item \textbf{What safeguards are in place against misuse?}
			\begin{itemize}
				\item Guardrails
				\item Human review triggers
				\item Usage limits
				\item Logs and monitoring
			\end{itemize}
		\end{enumerate}
		
		\vspace{0.1cm}
		\textbf{Why These Questions Matter:}
		\begin{itemize}
			\item Unpredictability can lead to harmful outcomes
			\item Safeguards prevent misuse and limit damage
			\item Monitoring enables detection and response
		\end{itemize}
		
		\vspace{0.1cm}
		\textbf{Key Point:}
		\begin{itemize}
			\item Technical controls are essential for safe deployment
			\item Test for edge cases and unexpected behaviors
		\end{itemize}
		
		\textcolor{myblue}{\textbf{Key Insight:}} Technical risk controls are \textcolor{myblue}{essential for safe GenAI deployment}!
	\end{frame}
	
	\begin{frame}{Slide 80: Organizational Questions}
		\fontsize{6.5}{7.5}\selectfont
		\textbf{Organizational Structure and Incentive Questions:}
		
		\vspace{0.1cm}
		\textbf{Key Questions:}
		\begin{enumerate}
			\item \textbf{Who owns risk for GenAI?}
			\begin{itemize}
				\item Is that person/team empowered?
				\item Clear accountability
			\end{itemize}
			\item \textbf{Are staff appropriately trained?}
			\begin{itemize}
				\item Know of unpredictable behavior
				\item Understand risks and limitations
				\item Prevent unintended consequences
			\end{itemize}
			\item \textbf{Are incentives aligned with responsible AI use?}
			\begin{itemize}
				\item Not just speed to deploy
				\item Risk-aware behaviors rewarded
			\end{itemize}
		\end{enumerate}
		
		\vspace{0.1cm}
		\textbf{Why These Questions Matter:}
		\begin{itemize}
			\item Ownership ensures accountability
			\item Training prevents misuse
			\item Incentives drive behavior
		\end{itemize}
		
		\textcolor{myblue}{\textbf{Key Insight:}} Organizational structure must \textcolor{myblue}{support responsible GenAI use}!
	\end{frame}
	
	\begin{frame}{Slide 81: External Environment Questions}
		\fontsize{6.5}{7.5}\selectfont
		\textbf{External Environment Questions:}
		
		\vspace{0.1cm}
		\textbf{Key Questions:}
		\begin{enumerate}
			\item \textbf{Are we prepared for changing laws and regulation?}
			\begin{itemize}
				\item EU AI Act
				\item US state legislation
				\item Compliance frameworks
			\end{itemize}
			\item \textbf{Are we aligned with best practices?}
			\begin{itemize}
				\item NIST AI RMF
				\item ISO/IEC 42001
				\item Industry standards
			\end{itemize}
		\end{enumerate}
		
		\vspace{0.1cm}
		\textbf{Why These Questions Matter:}
		\begin{itemize}
			\item Regulatory landscape is evolving
			\item Compliance is mandatory
			\item Best practices build trust
		\end{itemize}
		
		\vspace{0.1cm}
		\textbf{Key Point:}
		\begin{itemize}
			\item Following recognized frameworks strengthens compliance
			\item Builds trust with external stakeholders
			\item Demonstrates commitment to responsible AI
		\end{itemize}
		
		\textcolor{myblue}{\textbf{Key Insight:}} External alignment is \textcolor{myblue}{essential for regulatory compliance and stakeholder trust}!
	\end{frame}
	
	\begin{frame}{Slide 82: Questions for Risk Professionals Summary}
		\fontsize{6.5}{7.5}\selectfont
		\textbf{Questions for Risk Professionals - Summary:}
		
		\vspace{0.1cm}
		\textbf{Strategic Risk:}
		\begin{itemize}
			\item Is GenAI the right tool?
			\item What are the risks?
		\end{itemize}
		
		\vspace{0.1cm}
		\textbf{Model Selection:}
		\begin{itemize}
			\item What model is being used?
			\item Do we control it?
		\end{itemize}
		
		\vspace{0.1cm}
		\textbf{Data Governance:}
		\begin{itemize}
			\item What data is being used?
			\item Are protection policies in place?
		\end{itemize}
		
		\vspace{0.1cm}
		\textbf{Technical Controls:}
		\begin{itemize}
			\item Have we tested for unpredictability?
			\item What safeguards are in place?
		\end{itemize}
		
		\vspace{0.1cm}
		\textbf{Organizational:}
		\begin{itemize}
			\item Who owns risk?
			\item Are staff trained?
			\item Are incentives aligned?
		\end{itemize}
		
		\vspace{0.1cm}
		\textbf{External:}
		\begin{itemize}
			\item Are we prepared for regulation?
			\item Are we aligned with best practices?
		\end{itemize}
		
		\textcolor{myblue}{\textbf{Key Insight:}} A structured questioning approach \textcolor{myblue}{supports comprehensive GenAI governance}!
	\end{frame}
	
	\begin{frame}{Slide 83: Final Thoughts - Team Qualification}
		\fontsize{6.5}{7.5}\selectfont
		\textbf{Final Thoughts - Team Qualification:}
		
		\vspace{0.1cm}
		\textbf{Validation Team Requirements:}
		\begin{itemize}
			\item Knowledge of AI/ML techniques
			\item Hands-on expertise
			\item Strong foundations in:
			\begin{itemize}
				\item Product knowledge
				\item Econometrics
				\item Statistics
				\item Computer science
			\end{itemize}
			\item Proactive approach to staying abreast of latest advancements
		\end{itemize}
		
		\vspace{0.1cm}
		\textbf{Additional Considerations:}
		\begin{itemize}
			\item Augment validation staff
			\item Add AI/ML validation team
			\item Hands-on knowledge of:
			\begin{itemize}
				\item AI/ML development platforms
				\item Cloud-based environments
				\item Low-code and no-code environments
				\item CI/CD and MLOps platforms
			\end{itemize}
		\end{itemize}
		
		\vspace{0.1cm}
		\textbf{Key Point:}
		\begin{itemize}
			\item Familiarity with AI/ML technologies beyond traditional knowledge
			\item Specific training may be necessary
			\item Incorporate specialists
			\item Use external experts if needed
		\end{itemize}
		
		\textcolor{myblue}{\textbf{Key Insight:}} Validation teams need \textcolor{myblue}{updated skills and expertise for AI/ML models}!
	\end{frame}
	
	\begin{frame}{Slide 84: Final Thoughts - End-to-End Review}
		\fontsize{6.5}{7.5}\selectfont
		\textbf{Final Thoughts - End-to-End Review:}
		
		\vspace{0.1cm}
		\textbf{Review Framework:}
		\begin{itemize}
			\item Entire model validation framework needs review
			\item Regulators consider AI/ML models within scope of model risk management
			\item SR 11-7, being principle-based, applies to AI/ML models
		\end{itemize}
		
		\vspace{0.1cm}
		\textbf{Implications:}
		\begin{itemize}
			\item Challenges in model design
			\item Implications related to:
			\begin{itemize}
				\item Data
				\item Implementation
				\item Monitoring
				\item Documentation
				\item Use
			\end{itemize}
		\end{itemize}
		
		\vspace{0.1cm}
		\textbf{Key Point:}
		\begin{itemize}
			\item AI/ML models fall within existing regulatory frameworks
			\item Comprehensive review is necessary
		\end{itemize}
		
		\textcolor{myblue}{\textbf{Key Insight:}} End-to-end review is \textcolor{myblue}{essential for AI/ML model validation}!
	\end{frame}
	
	\begin{frame}{Slide 85: Final Thoughts - Balance and GenAI}
		\fontsize{6.5}{7.5}\selectfont
		\textbf{Final Thoughts - Balance and GenAI:}
		
		\vspace{0.1cm}
		\textbf{Model Complexity Balance:}
		\begin{itemize}
			\item More complex is not always better
			\item Balance model performance with:
			\begin{itemize}
				\item Interpretability
				\item Feedback from supervisors
				\item Cost and effort of implementation
				\item Maintenance and monitoring
				\item In-house expertise
				\item Availability of reputable code libraries
				\item Academic underpinning
			\end{itemize}
			\item Validation teams need balanced approach:
			\begin{itemize}
				\item Not just regulatory-driven
				\item Not just performance-oriented
				\item Nuanced cost-benefit analysis
			\end{itemize}
		\end{itemize}
		
		\vspace{0.1cm}
		\textbf{GenAI Governance:}
		\begin{itemize}
			\item Requires particular focus
			\item Support development of shared standards
			\item Transparency around GenAI
			\item Longer-term credibility depends on:
			\begin{itemize}
				\item Internal compliance
				\item Broader governance ecosystem
			\end{itemize}
		\end{itemize}
		
		\textcolor{myblue}{\textbf{Key Insight:}} Balance and GenAI governance are \textcolor{myblue}{essential for sustainable model risk management}!
	\end{frame}
	
	% ============================================================
	% SECTION 8: CHAPTER QUESTIONS - SLIDES 86-100
	% ============================================================
	
	\begin{frame}{Slide 86: Model Development and Testing - Key Points}
		\fontsize{6.5}{7.5}\selectfont
		\textbf{Key Points on Model Development and Testing:}
		
		\vspace{0.1cm}
		\textbf{Development Steps:}
		\begin{enumerate}
			\item Define objectives and scope
			\item Data collection and preprocessing
			\item Data analysis and visualization
			\item Feature engineering
			\item Model selection
			\item Model training/calibration
			\item Developer testing
			\item Documentation
		\end{enumerate}
		
		\vspace{0.1cm}
		\textbf{Testing Types:}
		\begin{itemize}
			\item Unit tests
			\item Component tests
			\item Integration tests
			\item System tests
			\item Performance tests
			\item Regression tests
			\item Acceptance tests
			\item Adversarial testing
		\end{itemize}
		
		\textcolor{myblue}{\textbf{Key Insight:}} Model development and testing are \textcolor{myblue}{systematic processes requiring comprehensive approaches}!
	\end{frame}
	
	\begin{frame}{Slide 87: Model Validation - Key Points}
		\fontsize{6.5}{7.5}\selectfont
		\textbf{Key Points on Model Validation:}
		
		\vspace{0.1cm}
		\textbf{Validation Lifecycle:}
		\begin{enumerate}
			\item Model identification
			\item Inventory
			\item Initial validation
			\item Production and monitoring
			\item Periodic validation
			\item Change-based validation
			\item Retirement
		\end{enumerate}
		
		\vspace{0.1cm}
		\textbf{Validation Activities:}
		\begin{itemize}
			\item Data assessment
			\item Assumptions and parameters
			\item Limitations identification
			\item Implementation review
			\item Conceptual soundness
			\item Performance monitoring
			\item Documentation review
			\item Controls review
			\item Outcomes analysis
			\item Testing and backtesting
		\end{itemize}
		
		\textcolor{myblue}{\textbf{Key Insight:}} Validation is a \textcolor{myblue}{comprehensive, ongoing process throughout the model lifecycle}!
	\end{frame}
	
	\begin{frame}{Slide 88: Model Misinterpretation - Key Points}
		\fontsize{6.5}{7.5}\selectfont
		\textbf{Key Points on Misinterpretation:}
		
		\vspace{0.1cm}
		\textbf{Common Sources:}
		\begin{enumerate}
			\item Lack of comprehension
			\item Neglecting model limitations
			\item Overemphasis on point estimates
			\item Failure to consider qualitative factors
			\item Ignoring decision context
			\item Confirmation bias
		\end{enumerate}
		
		\vspace{0.1cm}
		\textbf{Prevention Strategies:}
		\begin{itemize}
			\item Training and education
			\item Clear communication
			\item Open dialogue
			\item Supportive culture
		\end{itemize}
		
		\vspace{0.1cm}
		\textbf{Key Point:}
		\begin{itemize}
			\item Misinterpretation can have significant implications
			\item Requires active management
			\item Culture of critical thinking is essential
		\end{itemize}
		
		\textcolor{myblue}{\textbf{Key Insight:}} Misinterpretation can be prevented through \textcolor{myblue}{education, communication, and culture}!
	\end{frame}
	
	\begin{frame}{Slide 89: Three Lines of Defense - Key Points}
		\fontsize{6.5}{7.5}\selectfont
		\textbf{Three Lines of Defense - Key Points:}
		
		\vspace{0.1cm}
		\textbf{First Line (Model Owners/Developers):}
		\begin{itemize}
			\item Own the model and its performance
			\item Design and implement models
			\item Operate performance dashboards
			\item First to react to breaches
		\end{itemize}
		
		\vspace{0.1cm}
		\textbf{Second Line (Risk Management/Validation):}
		\begin{itemize}
			\item Set model risk management policy
			\item Conduct independent validations
			\item Challenge First Line assessments
			\item Maintain model inventory
		\end{itemize}
		
		\vspace{0.1cm}
		\textbf{Third Line (Internal Audit):}
		\begin{itemize}
			\item Provide independent assurance
			\item Review framework effectiveness
			\item Ensure policy compliance
		\end{itemize}
		
		\textcolor{myblue}{\textbf{Key Insight:}} Three lines provide \textcolor{myblue}{checks and balances in model risk management}!
	\end{frame}
	
	\begin{frame}{Slide 90: Model Changes - Key Points}
		\fontsize{6.5}{7.5}\selectfont
		\textbf{Model Changes - Key Points:}
		
		\vspace{0.1cm}
		\textbf{Change Management Process:}
		\begin{enumerate}
			\item Change requests
			\item Impact assessment
			\item Documentation
			\item Validation
			\item Approval
			\item Implementation
			\item Monitoring
		\end{enumerate}
		
		\vspace{0.1cm}
		\textbf{AI/ML Considerations:}
		\begin{itemize}
			\item More frequent monitoring needed
			\item Broader stakeholder involvement
			\item Focus on explainability
			\item Benchmarking against traditional models
			\item Assessment on different data sets
			\item Assessment of specific cases
			\item Back testing and cross-validation
			\item Reporting for bias detection
		\end{itemize}
		
		\textcolor{myblue}{\textbf{Key Insight:}} AI/ML model changes require \textcolor{myblue}{more frequent review and broader stakeholder involvement}!
	\end{frame}
	
	\begin{frame}{Slide 91: Data Governance - Key Points}
		\fontsize{6.5}{7.5}\selectfont
		\textbf{Data Governance - Key Points:}
		
		\vspace{0.1cm}
		\textbf{Core Components:}
		\begin{enumerate}
			\item \textbf{Strategy:}
			\begin{itemize}
				\item Vision and goals
				\item Policies and authorization
			\end{itemize}
			\item \textbf{Quality:}
			\begin{itemize}
				\item Accuracy, consistency, integrity
				\item Continuous monitoring
			\end{itemize}
			\item \textbf{Provenance:}
			\begin{itemize}
				\item Legal right to use data
				\item Copyright compliance
			\end{itemize}
			\item \textbf{Classification:}
			\begin{itemize}
				\item Structure and confidentiality
				\item Access controls
			\end{itemize}
			\item \textbf{Security:}
			\begin{itemize}
				\item Protection measures
				\item Regulatory compliance
			\end{itemize}
		\end{enumerate}
		
		\textcolor{myblue}{\textbf{Key Insight:}} Data governance is the \textcolor{myblue}{foundation for responsible AI and model management}!
	\end{frame}
	
	\begin{frame}{Slide 92: Model Inventory - Key Points}
		\fontsize{6.5}{7.5}\selectfont
		\textbf{Model Inventory - Key Points:}
		
		\vspace{0.1cm}
		\textbf{What to Include:}
		\begin{enumerate}
			\item Model identification (owner, user, developer)
			\item Purpose and description
			\item Methodology and classification
			\item Development source (in-house/vendor)
			\item Assumptions and limitations
			\item Documentation locations
			\item History of updates and validations
			\item Performance monitoring frequency
			\item Regulatory use indication
			\item Risk tier
		\end{enumerate}
		
		\vspace{0.1cm}
		\textbf{AI/ML Additional Considerations:}
		\begin{itemize}
			\item Complexity of methodology
			\item Data usage and quality
			\item Output parameters
			\item Model recalibration needs
			\item Risk ranking factors
			\item Performance monitoring metrics
		\end{itemize}
		
		\textcolor{myblue}{\textbf{Key Insight:}} Model inventory provides \textcolor{myblue}{comprehensive tracking and transparency} for all models!
	\end{frame}
	
	\begin{frame}{Slide 93: Model Risk Ranking - Key Points}
		\fontsize{6.5}{7.5}\selectfont
		\textbf{Model Risk Ranking - Key Points:}
		
		\vspace{0.1cm}
		\textbf{Key Factors:}
		\begin{enumerate}
			\item \textbf{Materiality:}
			\begin{itemize}
				\item Economic consequences
				\item Operational impact
				\item Reputational consequences
				\item Usage by different parties
				\item Impact on decisions, financial statements, regulatory reporting
			\end{itemize}
			\item \textbf{Complexity:}
			\begin{itemize}
				\item May elevate risk ranking
				\item Even if materiality is low
			\end{itemize}
			\item \textbf{Exposure:}
			\begin{itemize}
				\item Published financial reports
				\item Wide impact
				\item Regulatory significance
			\end{itemize}
		\end{enumerate}
		
		\vspace{0.1cm}
		\textbf{Important Note:}
		\begin{itemize}
			\item Minimum standards for all models regardless of materiality
			\item Absence of assessment introduces its own risk
		\end{itemize}
		
		\textcolor{myblue}{\textbf{Key Insight:}} Risk ranking determines \textcolor{myblue}{validation frequency and intensity}!
	\end{frame}
	
	\begin{frame}{Slide 94: Performance Monitoring - Key Points}
		\fontsize{6.5}{7.5}\selectfont
		\textbf{Performance Monitoring - Key Points:}
		
		\vspace{0.1cm}
		\textbf{Classification Models:}
		\begin{itemize}
			\item Accuracy, Precision, Recall, F1, AUC
			\item Confusion Matrices
		\end{itemize}
		
		\vspace{0.1cm}
		\textbf{Regression Models:}
		\begin{itemize}
			\item MAE, MSE, R-squared, RMSE
		\end{itemize}
		
		\vspace{0.1cm}
		\textbf{Additional Metrics:}
		\begin{itemize}
			\item Clustering metrics
			\item Deep learning metrics
			\item GenAI performance metrics
		\end{itemize}
		
		\vspace{0.1cm}
		\textbf{Key Point:}
		\begin{itemize}
			\item Performance monitoring must be tailored to model type
			\item Metrics must be appropriate for the model
			\item Model validation should assess metric appropriateness
		\end{itemize}
		
		\textcolor{myblue}{\textbf{Key Insight:}} Performance metrics must be \textcolor{myblue}{appropriate for the specific model type}!
	\end{frame}
	
	\begin{frame}{Slide 95: GenAI Validation - Key Points}
		\fontsize{6.5}{7.5}\selectfont
		\textbf{GenAI Validation - Key Points:}
		
		\vspace{0.1cm}
		\textbf{Special Considerations:}
		\begin{enumerate}
			\item Open-ended output
			\item Different performance metrics
			\item Data drift and coverage
			\item Hallucination
			\item Discrimination and bias
			\item Multi-modal issues
			\item Interpretability and explainability
			\item Rapid technological change
			\item Third-party and open-source models
			\item Computational resources
		\end{enumerate}
		
		\vspace{0.1cm}
		\textbf{Essential:}
		\begin{itemize}
			\item Human review is essential
			\item Regular validation and recalibration
			\item Particular attention to third-party models
		\end{itemize}
		
		\textcolor{myblue}{\textbf{Key Insight:}} GenAI validation requires \textcolor{myblue}{addressing unique challenges and human review}!
	\end{frame}
	
	\begin{frame}{Slide 96: GenAI Governance - Key Points}
		\fontsize{6.5}{7.5}\selectfont
		\textbf{GenAI Governance - Key Points:}
		
		\vspace{0.1cm}
		\textbf{Internal Levers:}
		\begin{enumerate}
			\item Clear ownership and accountability
			\item Risk-based review and deployment
			\item Incentivize responsible behavior
			\item Staff training
			\item Validation and monitoring
		\end{enumerate}
		
		\vspace{0.1cm}
		\textbf{External Constraints:}
		\begin{enumerate}
			\item Track legal and regulatory requirements
			\item Evaluate responsibilities with providers
			\item Align with emerging standards
		\end{enumerate}
		
		\vspace{0.1cm}
		\textbf{Adaptive Tools:}
		\begin{enumerate}
			\item Continuous monitoring
			\item Red-teaming and adversarial testing
			\item Scenario planning
			\item Technical safeguards
			\item Feedback loops
		\end{enumerate}
		
		\textcolor{myblue}{\textbf{Key Insight:}} GenAI governance requires \textcolor{myblue}{internal, external, and adaptive approaches}!
	\end{frame}
	
	\begin{frame}{Slide 97: Questions for Risk Professionals - Summary}
		\fontsize{6.5}{7.5}\selectfont
		\textbf{Questions for Risk Professionals - Summary:}
		
		\vspace{0.1cm}
		\textbf{Strategic Risk:}
		\begin{itemize}
			\item Is GenAI the right tool?
			\item What are the risks?
		\end{itemize}
		
		\vspace{0.1cm}
		\textbf{Model Selection:}
		\begin{itemize}
			\item What model is being used?
			\item Do we control it?
		\end{itemize}
		
		\vspace{0.1cm}
		\textbf{Data Governance:}
		\begin{itemize}
			\item What data is being used?
			\item Are protection policies in place?
		\end{itemize}
		
		\vspace{0.1cm}
		\textbf{Technical Controls:}
		\begin{itemize}
			\item Have we tested for unpredictability?
			\item What safeguards are in place?
		\end{itemize}
		
		\vspace{0.1cm}
		\textbf{Organizational:}
		\begin{itemize}
			\item Who owns risk?
			\item Are staff trained?
			\item Are incentives aligned?
		\end{itemize}
		
		\vspace{0.1cm}
		\textbf{External:}
		\begin{itemize}
			\item Are we prepared for regulation?
			\item Are we aligned with best practices?
		\end{itemize}
		
		\textcolor{myblue}{\textbf{Key Insight:}} Structured questioning \textcolor{myblue}{supports comprehensive GenAI governance}!
	\end{frame}
	
	\begin{frame}{Slide 98: Implementing AI Governance - BIS Survey}
		\fontsize{6.5}{7.5}\selectfont
		\textbf{Implementing AI Governance - BIS Survey Findings:}
		
		\vspace{0.1cm}
		\textbf{Data Governance:}
		\begin{itemize}
			\item Adequate documentation
			\item Data quality (accuracy, transparency, provenance)
			\item Security and responsible use
			\item Comprehensive data management systems
			\item Metadata on sources, version, curation
			\item Asset inventories and metadata registries
		\end{itemize}
		
		\vspace{0.1cm}
		\textbf{Organizational Structures:}
		\begin{itemize}
			\item Centralized solutions (steering committees, PMOs)
			\item Hub-and-spoke systems
			\item Decentralized approaches
			\item Accountability assignments
		\end{itemize}
		
		\textcolor{myblue}{\textbf{Key Insight:}} BIS survey shows \textcolor{myblue}{diverse approaches to AI governance implementation}!
	\end{frame}
	
	\begin{frame}{Slide 99: Implementing AI Governance (Continued)}
		\fontsize{6.5}{7.5}\selectfont
		\textbf{Implementing AI Governance - BIS Survey (Continued):}
		
		\vspace{0.1cm}
		\textbf{Guidelines and Policies:}
		\begin{itemize}
			\item Terms of use for AI systems
			\item Guidelines for responsible AI use
			\item Privacy standards
			\item Explainability requirements
			\item International principles
			\item Documentation of assessment methods
		\end{itemize}
		
		\vspace{0.1cm}
		\textbf{Risk Management Frameworks:}
		\begin{itemize}
			\item Regular audit of AI systems
			\item Evaluation of safety measures
			\item Incident response procedures
			\item Business continuity plans
			\item Skills maintenance processes
			\item Leverage existing frameworks (3LoD)
		\end{itemize}
		
		\textcolor{myblue}{\textbf{Key Insight:}} Comprehensive guidelines and risk management are \textcolor{myblue}{essential for AI governance}!
	\end{frame}
	
	\begin{frame}{Slide 100: Data and AI Model Governance - Final Summary}
		\fontsize{6.5}{7.5}\selectfont
		\textbf{Data and AI Model Governance - Comprehensive Summary:}
		
		\vspace{0.1cm}
		\textbf{Key Pillars:}
		\begin{enumerate}
			\item \textbf{Data Governance:}
			\begin{itemize}
				\item Strategy, quality, provenance, classification, security
			\end{itemize}
			\item \textbf{Model Governance:}
			\begin{itemize}
				\item Policies, documentation, inventory, risk ranking
			\end{itemize}
			\item \textbf{Model Development:}
			\begin{itemize}
				\item Objectives, data, features, training, testing
			\end{itemize}
			\item \textbf{Model Validation:}
			\begin{itemize}
				\item Initial, ongoing, change-based
				\item Comprehensive review activities
			\end{itemize}
			\item \textbf{Three Lines of Defense:}
			\begin{itemize}
				\item Owners/developers, risk management, internal audit
			\end{itemize}
			\item \textbf{GenAI Governance:}
			\begin{itemize}
				\item Internal levers, external constraints, adaptive tools
			\end{itemize}
		\end{enumerate}
		
		\textcolor{myred}{\textbf{Exam Success:}} Understand \textcolor{myblue}{data governance, model risk management, validation, and GenAI governance}!
	\end{frame}
	
\end{document}